\documentclass[11pt,a4paper]{article}
\usepackage[T1]{fontenc}
\usepackage[utf8]{inputenc}
\usepackage{lmodern}
\usepackage[a4paper,margin=26mm,headheight=14pt]{geometry}
\usepackage{amsmath,amssymb,amsthm,mathtools,mathrsfs}
\usepackage{booktabs,array,longtable}
\usepackage{microtype}
\usepackage{enumitem}
\usepackage{xcolor}
\usepackage{fancyhdr}
\usepackage{hyperref}
\definecolor{linkblue}{RGB}{24,63,105}
\hypersetup{colorlinks=true,linkcolor=linkblue,citecolor=linkblue,urlcolor=linkblue,
 pdftitle={Joint Polygamma Moments, Gap-Eight Reconstruction, and Multiscale Stieltjes Resonances},
 pdfauthor={Research continuation for the ProveIt project},
 pdfsubject={Exact special-function identities and rigorous research continuation}}
\urlstyle{same}
\Urlmuskip=0mu plus 2mu\relax
\setlength{\emergencystretch}{2em}
\setlength{\parskip}{3pt}
\setlist[enumerate]{leftmargin=*,itemsep=5pt,topsep=5pt}
\setlist[itemize]{leftmargin=*,itemsep=3pt,topsep=4pt}
\numberwithin{equation}{section}
\newtheorem{theorem}{Theorem}[section]
\newtheorem{lemma}[theorem]{Lemma}
\newtheorem{proposition}[theorem]{Proposition}
\newtheorem{corollary}[theorem]{Corollary}
\theoremstyle{definition}
\newtheorem{definition}[theorem]{Definition}
\theoremstyle{remark}
\newtheorem{remark}[theorem]{Remark}
\pagestyle{fancy}
\fancyhf{}
\fancyhead[L]{\small ProveIt: special-function research continuation}
\fancyhead[R]{\small 12 October 2026}
\fancyfoot[C]{\thepage}
\renewcommand{\headrulewidth}{0.3pt}
\title{\bfseries Joint Polygamma Moments,\\
Gap-Eight Reconstruction, and\\
Multiscale Stieltjes Resonances}
\author{Research continuation for the ProveIt project}
\date{12 October 2026}

\begin{document}
\maketitle
\begin{abstract}
We advance three explicitly open directions in the current ProveIt
polylogarithm and special-function reports. First, an unpaired deformation
of a symplectic Gamma integral separates every joint polygamma cumulant
whose total derivative weight is at most $2n+2$. Positive-order statistics
obey an exact finite Wick rule, whereas the digamma sector has explicitly
nonzero higher cumulants. A small-parameter limit proves that the Gaussian
range cannot be extended universally. Second, we make fixed-gap harmonic
tail reconstruction effective for every prescribed even gap and complete
gap eight for all indices. The proof combines explicit Gamma-integral
asymptotics with an exact finite-field certificate for the remaining
$8187$ matrices; it also produces an all-order family of ordinary,
absolutely convergent polygamma-product sums equal to $-1/2$.
Third, we evaluate the constant term in collisions of arbitrary argument
derivatives of generalized Stieltjes constants at two positive monomial
scales. Three elementary resonance equations describe all new contacts.
Their coefficients reduce to zeta derivatives, harmonic numbers, and
Stieltjes constants at one. A first-index specialization realizes every
integer trigamma tail as a collision contact, with a polylogarithmic
generator and explicit primitive. Barnes-integral companion evaluations,
source qualifications, reproducible verification programs, and a precise
research agenda accompany the proofs. Novelty is asserted relative to
the inspected repository baseline, with classical ingredients identified.
\end{abstract}
\tableofcontents
\clearpage

% BEGIN sections/00_introduction.tex
\section{Scope, principal results, and conventions}
\label{intro:section}

The starting point is the ProveIt repository at revision
\texttt{03112a7d3b91} (full revision in \cite{ProveIt}), inspected on
12 October 2026. We read the polylogarithm manuscript and the text sources
in the incoming reports available at that revision; the most immediate
predecessors are the ordered-boundary, gap-six, and Gamma-normalization
reports \cite{ProveIt,OrderedIncoming,GapSix,GammaIncoming}.
This matters because several questions suggested by older manuscript
passages have already been answered in newer incoming material.
In particular, complete gap six, fixed-rate Stieltjes collisions, and
the paired Gamma deformation are prior results, not new contributions
of this article.

Our aim is to obtain exact evaluations and structural formulas. Estimates
are used where needed to prove identities or to turn eventual
nonsingularity into a theorem valid for all indices. They are not the
principal conclusions. Each of the three main directions gives a
finite, explicit way to evaluate an infinite family.

\subsection{Three advances and their exact scope}

\paragraph{Joint polygamma moments.}
For the positive $n$-variable Gamma density defined in
\eqref{gc:density}, introduce
$R_q=\sum_{j=1}^n\Re\psi^{(q)}(\sigma+ix_j/2)$.
Theorem~\ref{gc:main} evaluates all joint cumulants satisfying
$\sum_i(q_i+1)\le2n+2$. It gives means, covariances, and the particularly
simple higher-cumulant formula
\[
 \kappa(R_{q_1},\ldots,R_{q_r})
 =-2^{-r}\psi^{(r-1)}((2n+2)\sigma)
          \prod_i\delta_{q_i0},\qquad r\ge3.
\]
Thus the positive-order sector has exact Gaussian moments in a finite
weighted range, while the digamma sector retains a non-Gaussian
correction. This answers the joint-separation question of the latest
Gamma report in that range. The partition formula
\eqref{gc:moment} evaluates the corresponding ordinary integrals without
leaving parameter derivatives unevaluated. Proposition~\ref{gc:sharpness}
also proves a specific obstruction to extending the source's
single-trigamma matching degree universally.

\paragraph{Gap-eight reconstruction and sums.}
For $p=2b-2m$, $q=2b$, the harmonic-tail problem reduces to a rational
matrix $M_{b,m}$. Theorem~\ref{eg:effective} gives a fully computable
threshold $B(m)$ beyond which it is nonsingular, for every prescribed
$m\ge1$. Theorem~\ref{eg:eight-main} completes $m=4$ at every $b\ge5$,
settling the explicitly open gap-eight ray. Its finite step is an exact
computation, not a numerical extrapolation. Corollary~\ref{eg:eight-identities}
then gives rational polynomials $P_b,S_b$ such that
\[
 \sum_{n=0}^{\infty}\left[
 P_b(n+\tfrac12)\zeta(2b-8,n+1)\zeta(2b,n+1)
       -S_b(n+\tfrac12)\right]=-\frac12.
\]
The subtracted series is absolutely convergent. Matrix rank is taken on
the indicated formal coordinates; it is not a claim that the evaluated
zeta values are rationally or algebraically independent.

\paragraph{Different-power Stieltjes collisions.}
Theorem~\ref{mc:main} evaluates the change in every finite-part constant
when the two arguments are shifted by
$a=\lambda\epsilon^\alpha$, $b=\mu\epsilon^\beta$, with positive
parameters and unequal exponents. The equal-exponent case is inherited.
The answer consists of two explicit mixed contacts and at most one
singular--singular contact. The latter can survive precisely when
\[
 j(\alpha-\beta)=\beta(r+t+1),\qquad j\in\mathbb Z_{\ge1},
 \quad\alpha>\beta.
\]
The mixed contacts obey two further integer selection rules. This
classification explains constants absent from the fixed-rate theorem,
including the two renormalized limits $-1-\zeta(3)$ and $-1+\zeta(4)$
in \eqref{mc:q2limit}--\eqref{mc:q3limit}.
An especially short consequence is
\[
 \operatorname{CT}_{\epsilon}\int_0^1
 \gamma_1(x+\epsilon^{j+1})\gamma_0(x+\epsilon^j)\,dx
 =\frac32\gamma_2+2\gamma\zeta(2)-\zeta(3)+\psi'(j+1).
\]
Here the fixed finite part is inherited; the additional trigamma tail
is the newly identified collision term.

\subsection{Classical ingredients and the novelty boundary}

The multidimensional Gamma integral is a classical Gustafson integral;
we use the explicit symplectic normalization reproduced in
\cite[Eq.~(2)]{DerkachovManashov}, together with the repository's
determinant transformation. The new step is separation in independent
power-sum coordinates through their exact injectivity range.
The reconstruction proof builds on the repository's Bernoulli recurrence
and Gamma basis. Its new steps are a uniform effective threshold,
the complete gap-eight certificate, and the associated evaluated sums.

Finite-part integration and generalized Stieltjes-transform asymptotics
have an established literature \cite{TicaGalapon}; integrals of products
of Hurwitz zeta functions also occur in \cite{ShpotParis}. The monomial
collision theorem below uses elementary local subtraction, an incomplete
beta expansion, and cancellation of two meromorphic branches. Its new
content is the evaluated all-index contact law for the specified
two-scale problem. The Barnes companion in Section~\ref{bim:section}
is explicitly presented as a consequence of the classical first Barnes
lemma \cite[5.13.3]{DLMF}. We do not claim that parameter differentiation,
Bell polynomials, Barnes integration, or the resulting elementary
Euler-sum generators are new methods.

The comparison establishes advances over the inspected repository,
not an exhaustive historical-priority determination. The claims are
proved in the usual mathematical sense, with an explicitly isolated
exact computational component for the finite gap-eight interval.
The accompanying numerical calculations are independent diagnostics;
they are not proofs by integer-relation recognition. No proof-assistant
formalization is claimed.

\subsection{Notation and finite-part convention}

We use $\psi=\Gamma'/\Gamma$, $\psi^{(0)}=\psi$, and the standard
Hurwitz zeta function $\zeta(s,x)$. A superscript in parentheses on
$\gamma_m(x)$ denotes an argument derivative, whereas $m$ is its
Stieltjes index. The convention for generalized Stieltjes constants is
\[
 \zeta(1+u,x)=\frac1u+
       \sum_{m\ge0}\frac{(-1)^m}{m!}\gamma_m(x)u^m,
 \qquad \gamma_m=\gamma_m(1),\quad\gamma_0=\gamma.
\]
Thus $\gamma_0(x)=-\psi(x)$. Rising factorials are $(a)_k$;
falling factorials are explicitly underlined. We set
\[
 H_j^{(h)}=\sum_{k=1}^j k^{-h},\qquad H_j=H_j^{(1)},
 \qquad H_0^{(h)}=0.
\]
Bernoulli numbers have generating function
$t/(e^t-1)=\sum_{k\ge0}B_k t^k/k!$.
All logarithms of positive arguments are real; analytic logarithms
elsewhere are fixed locally at the stated base point.

The operation $\operatorname{CT}_{\epsilon}$ extracts the coefficient
of $\epsilon^0(\log\epsilon)^0$ from the specified power--logarithm
expansion. It is a coordinate-dependent finite part, not an unspecified
regularization and not necessarily a limit. We supply explicit ordinary
subtractions when stating a renormalized limit. The monomial exponents
are fixed while $\epsilon\downarrow0$; a jointly varying exponent ratio
near a resonance is a separate problem.

% END sections/00_introduction.tex

% BEGIN sections/01_joint_polygamma.tex
% Independent continuation, 12 October 2026.
% Fragment uses amsmath, amssymb, amsthm; theorem/lemma/corollary environments.
% Primary bibliography keys below are suggestions for the parent article.
\section{Joint Gamma--polygamma cumulants}
\label{gc:section}

The latest incoming Gamma-normalization report \cite{GammaIncoming} evaluates the appropriate
Gustafson determinant, establishes a paired parameter deformation, and
proves Gaussian moments for a single trigamma statistic through degree
$n+1$. Its research question Q2 asks for a precise classification of
joint polygamma moments. The theorem below answers this question in a
complete weighted range. It also includes the digamma statistic and all
even polygamma orders, which require the unpaired deformation.
The classical integral and the source's density are inherited; the
explicit cumulant classification and its consequences are the present
extension. No claim of priority over all literature is made.

For $n\ge1$, $\sigma>0$, put $N=2n+2$ and consider the probability
measure on $(0,\infty)^n$ with density proportional to
\begin{equation}\label{gc:density}
 \Delta(x_1^2,\ldots,x_n^2)
 \Delta(\cosh\pi x_1,\ldots,\cosh\pi x_n)
 \prod_{j=1}^n x_j\sinh(\pi x_j)
       |\Gamma(\sigma+ix_j/2)|^{2N},
\end{equation}
where $\Delta(y)=\prod_{j<k}(y_k-y_j)$. Its normalizer is
\begin{equation}\label{gc:normalizer}
 Z_{n,\sigma}
 =2^{n(3n+1)/2}\pi^{n(n+1)}n!
       \frac{\Gamma(2\sigma)^{(n+1)(2n+1)}}{\Gamma(N\sigma)}.
\end{equation}
The Vandermondes have the same sign, so this is a positive measure.
Write $\mathbb E=\mathbb E_{n,\sigma}$ and define, for $q\ge0$,
\begin{equation}\label{gc:R}
 R_q=\sum_{j=1}^n\Re\psi^{(q)}(\sigma+ix_j/2),
 \qquad \psi^{(0)}=\psi.
\end{equation}
The \emph{weight} of $R_q$ is $q+1$. A joint cumulant means
\[
 \kappa(Y_1,\ldots,Y_r)
 =\left.\partial_{t_1}\cdots\partial_{t_r}
     \log\mathbb E e^{t_1Y_1+\cdots+t_rY_r}\right|_{\mathbf t=0}.
\]
All fixed moments exist. Indeed the positive-order polygammas are bounded
on each vertical line $\Re z=\sigma$, the digamma grows at most
logarithmically, and the Gamma weight has an exponential margin after
multiplication by every fixed polynomial or logarithmic power.

\begin{theorem}[Explicit joint cumulants in the complete weighted range]
\label{gc:main}
Let $q_1,\ldots,q_r\ge0$ and suppose
\begin{equation}\label{gc:range}
                   \sum_{i=1}^r(q_i+1)\le N=2n+2.
\end{equation}
The one-point and two-point cumulants are
\begin{align}
 \mathbb E R_q
   &=\frac12\left[(N-2^q)\psi^{(q)}(2\sigma)
               -\delta_{q0}\psi(N\sigma)\right],
              \label{gc:mean}\\
 \operatorname{Cov}(R_p,R_q)
   &=\frac14\left[\psi^{(p+q+1)}(2\sigma)
               -\delta_{p0}\delta_{q0}\psi'(N\sigma)\right].
              \label{gc:cov}
\end{align}
For every $r\ge3$ in the same weighted range,
\begin{equation}\label{gc:higher}
 \boxed{\quad
 \kappa(R_{q_1},\ldots,R_{q_r})
 =-\frac{\psi^{(r-1)}(N\sigma)}{2^r}
                    \prod_{i=1}^r\delta_{q_i0}.
 \quad}
\end{equation}
Thus every mixed cumulant of order at least three that contains a
positive-order polygamma statistic vanishes in this range.
\end{theorem}

\subsection{Proof by independent parameter shifts}

For independent parameters $z_1,\ldots,z_N$ near zero, deform all Gamma
parameters to $\sigma+z_\ell$. Gustafson's symplectic Barnes integral,
in the normalization of Derkachov--Manashov \cite[Eq.~(2)]{DerkachovManashov}, and the
source's determinant-to-density transformation give
\begin{equation}\label{gc:unpaired}
\begin{split}
 &\mathbb E\prod_{j=1}^n\prod_{\ell=1}^N
 \frac{\Gamma(\sigma+z_\ell+ix_j/2)
       \Gamma(\sigma+z_\ell-ix_j/2)}
      {\Gamma(\sigma+ix_j/2)\Gamma(\sigma-ix_j/2)}\\
 &\qquad=
 \frac{\Gamma(N\sigma)}{\Gamma(N\sigma+\sum_\ell z_\ell)}
 \prod_{\ell<h}
       \frac{\Gamma(2\sigma+z_\ell+z_h)}{\Gamma(2\sigma)}.
\end{split}
\end{equation}
This is an identity of holomorphic germs. On a sufficiently small
polydisc every shifted Gamma parameter remains in the right half-plane.
Uniform Stirling bounds leave an integrable exponential majorant, even
after any fixed mixed derivative; they justify coefficient extraction
under the integral. Holomorphic log-Gamma is used on that half-plane,
and the analytic logarithm of the expectation is fixed by its value
zero at the origin.

Put $p_k=\sum_{\ell=1}^N z_\ell^k$. Taylor expansion of the logarithm
of the integrand in \eqref{gc:unpaired} gives
\[
 \log\left(\prod_{j,\ell}\frac{\Gamma(\sigma+z_\ell+ix_j/2)
     \Gamma(\sigma+z_\ell-ix_j/2)}
     {\Gamma(\sigma+ix_j/2)\Gamma(\sigma-ix_j/2)}\right)
            =2\sum_{k\ge1}\frac{R_{k-1}}{k!}p_k.
\]
For every integer $k\ge1$ the elementary pair-sum identity is
\begin{equation}\label{gc:pairsum}
 \sum_{\ell<h}(z_\ell+z_h)^k
 =(N-2^{k-1})p_k+\frac12\sum_{a=1}^{k-1}
                              \binom ka p_a p_{k-a}.
\end{equation}
The coefficient of $p_k$ follows by excluding the diagonal from
$\sum_{\ell,h}z_\ell^a z_h^{k-a}$; the two endpoint terms each
contribute $(N-1)p_k$ before the factor $1/2$ is applied.
Taking logarithms in \eqref{gc:unpaired} therefore gives
\begin{equation}\label{gc:power-sum-log}
\begin{split}
 &\log\mathbb E\exp\left(2\sum_{k\ge1}\frac{R_{k-1}}{k!}p_k\right)\\
 &\quad=\sum_{k\ge1}
       \frac{(N-2^{k-1})\psi^{(k-1)}(2\sigma)}{k!}p_k
       +\frac12\sum_{a,b\ge1}
          \frac{\psi^{(a+b-1)}(2\sigma)}{a!b!}p_a p_b
       -\sum_{r\ge1}\frac{\psi^{(r-1)}(N\sigma)}{r!}p_1^r.
\end{split}
\end{equation}

The next finite algebra observation is the step that permits separation
of the different statistics, rather than merely recording a generating
identity.

\begin{lemma}[Independent power-sum coordinates through degree $N$]
\label{gc:psindependent}
Assign weight $k$ to an indeterminate $P_k$. Substitution
$P_k\mapsto p_k(z_1,\ldots,z_N)$ is injective on the space of
polynomials of weighted degree at most $N$.
\end{lemma}
\begin{proof}
Only $P_1,\ldots,P_N$ can occur in that degree range.
Newton's identities give an invertible triangular change between
$p_1,\ldots,p_N$ and the first $N$ elementary symmetric polynomials,
which are algebraically independent. Alternatively the Jacobian of
the power sums is
\[
 \det\left(\frac{\partial p_k}{\partial z_\ell}\right)_{k,\ell=1}^N
              =N!\prod_{\ell<h}(z_h-z_\ell),
\]
which is not the zero polynomial. Either proof gives the claimed
injectivity.
\end{proof}

\begin{proof}[Proof of Theorem~\ref{gc:main}]
The homogeneous component of degree $d\le N$ in
\eqref{gc:power-sum-log} involves only $p_1,\ldots,p_d$.
By Lemma~\ref{gc:psindependent}, compare its coefficients as polynomials
in independent power-sum variables. Set $p_k=k!t_{k-1}/2$ in this
finite-degree comparison. The logarithmic moment generating function
then has, through weighted degree $N$, the form
\begin{equation}\label{gc:cgf}
\begin{split}
 \log\mathbb E\exp\left(\sum_{q\ge0}t_qR_q\right)
  ={}&\frac12\sum_{q\ge0}(N-2^q)\psi^{(q)}(2\sigma)t_q\\
    &+\frac18\sum_{p,q\ge0}\psi^{(p+q+1)}(2\sigma)t_pt_q
      -\sum_{r\ge1}\frac{\psi^{(r-1)}(N\sigma)}{2^r r!}t_0^r.
\end{split}
\end{equation}
Here equality means coefficient equality only through that weighted
degree, not an unrestricted identity of generating functions.
The linear, quadratic and higher coefficients give
\eqref{gc:mean}--\eqref{gc:higher}.
\end{proof}

\subsection{A finite formula for every retained ordinary moment}

Let $\mathcal P_r$ be the set of partitions of $\{1,\ldots,r\}$.
For a block $B$ define $c(B)$ by \eqref{gc:mean} when $|B|=1$,
by \eqref{gc:cov} when $|B|=2$, and by
\eqref{gc:higher} when $|B|\ge3$, using the orders attached to the
labels in $B$. Since every block satisfies \eqref{gc:range} whenever
the full list does, the ordinary moment is explicitly
\begin{equation}\label{gc:moment}
 \boxed{\qquad
 \mathbb E\prod_{i=1}^rR_{q_i}
             =\sum_{\pi\in\mathcal P_r}\prod_{B\in\pi}c(B),
 \qquad \sum_i(q_i+1)\le N.
 \qquad}
\end{equation}
This finite partition formula evaluates all the indicated integrals
using only polygamma values at $2\sigma$ and $N\sigma$. It contains no
unevaluated spectral derivative or auxiliary integral.

For positive orders define the centered statistics
\[
 Y_q=R_q-\frac12(N-2^q)\psi^{(q)}(2\sigma),\qquad q\ge1.
\]
All blocks of size other than two disappear from their centered moments.
Consequently
\begin{corollary}[Weighted Wick rule]
\label{gc:wick}
For $q_i\ge1$ and $\sum_i(q_i+1)\le N$,
\begin{equation}\label{gc:wickeq}
 \mathbb E\prod_{i=1}^rY_{q_i}
 =\begin{cases}
 0,&r\text{ odd},\\[2pt]
 \displaystyle 4^{-r/2}\sum_{\text{pairings }\pi}
         \prod_{\{i,j\}\in\pi}\psi^{(q_i+q_j+1)}(2\sigma),
                                    &r\text{ even}.
 \end{cases}
\end{equation}
\end{corollary}
For example, for $n\ge5$ the genuinely mixed fourth moment is
\begin{equation}\label{gc:example13}
 \mathbb E(Y_1^2Y_3^2)
 =\frac{\psi'''(2\sigma)\psi^{(7)}(2\sigma)}{16}
                +\frac{\psi^{(5)}(2\sigma)^2}{8}.
\end{equation}
At $\sigma=1/4$ this is $33\pi^{12}/2$, using
$\psi'''(1/2)=\pi^4$, $\psi^{(5)}(1/2)=8\pi^6$, and
$\psi^{(7)}(1/2)=136\pi^8$.
For $n\ge2$ another mixed evaluation is
\begin{equation}\label{gc:example012}
 \mathbb E\bigl[(R_0-\mathbb ER_0)(R_1-\mathbb ER_1)
                                  (R_2-\mathbb ER_2)\bigr]=0,
\end{equation}
whose total weight is $1+2+3=6$.
The odd-order restriction, $S_{2a-1}=R_{2a-1}$, reads
\[
 \mathbb ES_{2a-1}=(n+1-4^{a-1})\psi^{(2a-1)}(2\sigma),\qquad
 \operatorname{Cov}(S_{2a-1},S_{2b-1})
               =\tfrac14\psi^{(2a+2b-1)}(2\sigma)
\]
whenever, respectively, $a\le n+1$ and $a+b\le n+1$.
Every higher cumulant vanishes when the sum of the corresponding
$a$-indices is at most $n+1$. This is the requested separation of the
joint odd-polygamma coordinates.

\subsection{New scalar integral evaluations}

Set $n=1$, $\sigma=1/4$ and
\[
 \omega(x)=x\sinh(\pi x)|\Gamma(1/4+ix/2)|^8,
 \quad z_x=1/4+ix/2,
\]
\[
 g(x)=\Re\psi(z_x)+\gamma+3\log2,\qquad
 h(x)=\Re\psi'(z_x)-\pi^2/2.
\]
Here $\int_0^\infty\omega(x)\,dx=4\pi^5$.
The theorem, whose weighted budget is $N=4$, gives the exact ordinary
integrals, using the standard special values in \cite[\S5.15]{DLMF},
\begin{align}
 \int_0^\infty\omega(x)\Re\psi''(z_x)\,dx&=0,
                                                    \label{gc:scalar1}\\
 \int_0^\infty\omega(x)g(x)\Re\psi''(z_x)\,dx&=\pi^9,
                                                    \label{gc:scalar2}\\
 \int_0^\infty\omega(x)g(x)^2h(x)\,dx&=0,
                                                    \label{gc:scalar3}\\
 \int_0^\infty\omega(x)g(x)^3\,dx&=\pi^5\zeta(3),
                                                    \label{gc:scalar4}\\
 \int_0^\infty\omega(x)g(x)^4\,dx&=\frac{\pi^9}{15}.
                                                    \label{gc:scalar5}
\end{align}
For the last two, the centered digamma cumulants are
\[
 \kappa_2(R_0)=\frac{\pi^2}{12},\quad
 \kappa_3(R_0)=\frac{\zeta(3)}4,\quad
 \kappa_4(R_0)=-\frac{\pi^4}{240}.
\]
Thus the fourth central moment is
$3(\pi^2/12)^2-\pi^4/240=\pi^4/60$; multiplication by
$4\pi^5$ gives \eqref{gc:scalar5}.
The nonzero digamma cumulants explain why the Gaussian description
belongs to the positive-order sector.

\subsection{The finite range is essential}

At degree $N+1$, Newton's relation $e_{N+1}=0$ in $N$ variables creates
a nonzero kernel:
\begin{equation}\label{gc:firstrelation}
 \sum_{\lambda\vdash N+1}
       \frac{(-1)^{N+1-\ell(\lambda)}}{z_\lambda}
                 p_{\lambda_1}\cdots p_{\lambda_{\ell(\lambda)}}=0,
 \qquad z_\lambda=\prod_{j\ge1}j^{m_j}m_j!.
\end{equation}
One therefore cannot compare all formal power-sum coefficients beyond
the theorem's range. It is an algebraic failure of this extraction
argument, not a proof that every first excluded cumulant is nonzero.
The following limiting example supplies an actual obstruction to
extending the source's single-trigamma Gaussian law universally.

\begin{proposition}[A nonzero first excluded Gaussian moment]
\label{gc:sharpness}
For $n=1$ define the source-normalized statistic
\[
 X_{1,\sigma}=\frac{2R_1-2\psi'(2\sigma)}
                         {\sqrt{\psi'''(2\sigma)}}.
\]
Then
\begin{equation}\label{gc:skewlimit}
        \lim_{\sigma\downarrow0}\mathbb E X_{1,\sigma}^3
                       =\frac{2\sqrt6}{9}>0.
\end{equation}
In particular the Gaussian matching degree $n+1$ cannot be increased
to $n+2$ uniformly over all $n\ge1$ and all $\sigma>0$.
This proposition does not evaluate the general $(n+2)$nd defect.
\end{proposition}
\begin{proof}
At $n=1$, \eqref{gc:normalizer} is
$Z_{1,\sigma}=4\pi^2\Gamma(2\sigma)^6/\Gamma(4\sigma)$,
and $\sigma^5 Z_{1,\sigma}\to\pi^2/4$.
Set $x=2\sigma y$. Using $z\Gamma(z)\to1$ as $z\to0$, the
normalized density in the $y$ variable tends pointwise to
\[
                q_0(y)=\frac{32}{\pi}\frac{y^2}{(1+y^2)^4},
                \qquad y>0.
\]
Its integral is one by the beta integral. Since the approximating
densities also integrate to one, Scheff\'e's lemma gives convergence
in $L^1$. Also
\[
 \sigma^2\Re\psi'(\sigma(1+iy))
       \longrightarrow\Re(1+iy)^{-2},\qquad
 \sigma^2|\psi'(\sigma(1+iy))|
       \le\sigma^2\psi'(\sigma),
\]
and the upper bound tends to one. The density convergence and this
uniform bound give convergence of every fixed moment of the scaled
statistic. If $t=y^2/(1+y^2)$ then $t$ has the
$\operatorname{Beta}(3/2,5/2)$ law and
\[
 \Re(1+iy)^{-2}=f(t)=1-3t+2t^2.
\]
The beta moment formula
$\mathbb Et^j=(3/2)_j/(4)_j$ gives
\[
 \mathbb Ef=\tfrac14,\qquad
 \mathbb E(f-1/4)^2=\tfrac3{32},\qquad
 \mathbb E(f-1/4)^3=\tfrac1{64}.
\]
Finally $\sigma^2\psi'(2\sigma)\to1/4$ and
$\sigma^4\psi'''(2\sigma)\to3/8$. Thus the normalized variable
converges in moments to $\sqrt{32/3}(f-1/4)$ and its third moment
is $(32/3)^{3/2}/64=2\sqrt6/9$.
\end{proof}

\subsection{A joint Gaussian consequence}

Although the main conclusions are exact integral evaluations, the
finite moment law also has a useful joint limit. Fix distinct positive
orders $q_1,\ldots,q_d$ and define
\[
 C(\sigma)_{ij}=\tfrac14\psi^{(q_i+q_j+1)}(2\sigma).
\]
This matrix is positive definite: the polygamma integral representation
gives
\[
 v^TC(\sigma)v
 =\frac14\int_0^\infty
       \left(\sum_i(-1)^{q_i}v_i t^{q_i}\right)^2
                    \frac{t e^{-2\sigma t}}{1-e^{-t}}\,dt>0
\]
for every nonzero real $v$. Distinct monomial degrees ensure strictness.
For any positive sequence $\sigma_n$, let
\[
 \mathbf Z_n=C(\sigma_n)^{-1/2}
       (R_{q_i}-\mathbb ER_{q_i})_{i=1}^d.
\]
Every mixed moment of total ordinary degree $k$ is exactly the
corresponding standard $d$-variate Gaussian moment once
$k(1+\max_iq_i)\le2n+2$. This follows by expanding the deterministic
linear change of coordinates and applying the weighted Wick rule.
In particular $\mathbf Z_n$ converges in distribution to $N(0,I_d)$.
For a proof, every fixed linear projection has eventually the exact
Gaussian moments of each prescribed order; the second moments give
tightness and higher even moments give uniform integrability.
Moment determinacy of the normal law and the Cram\'er--Wold theorem
finish the argument. No estimate uniform in $\sigma_n$ is needed.

\subsection{Independent verification}

The accompanying program \texttt{verify\_gamma\_cumulants.py} performs
122 exact assertions. These verify the pair-sum identity at twelve
orders and eight numbers of variables, verify the ranks of the
power-sum substitution through the first excluded degree for up to five
variables, and independently evaluate the beta moments and displayed
quarter-parameter arithmetic. In particular the rank drops by one at
weighted degree $N+1$, as predicted by \eqref{gc:firstrelation}.

The numerical side uses continuous Cauchy--Binet in the opposite
direction: it integrates one-dimensional moment entries with inserted
products of polygammas, constructs a determinant of their finite Taylor
jets, and extracts moments and cumulants by finite polynomial
arithmetic. It does not use the Gamma quotient for the expectation
side. Seventeen checks include the zero mean of $R_2$, the third and
fourth centered digamma cumulants, mixed vanishing cumulants, and
$\operatorname{Var}(R_2)=\operatorname{Cov}(R_1,R_3)$ at $n=2$.
The cases use quarter and non-quarter shifts. At 36 decimal digits,
with 56 Gauss--Legendre nodes on each stated integration interval, the
largest scaled residual is approximately $8.2\times10^{-34}$.
A repeat at 46 decimal digits and 72 nodes per interval
reduces the largest scaled residual to approximately $2.1\times10^{-43}$.
A separate direct adaptive integral gives
\[
 \begin{aligned}
 \mathbb EX_{1,0.02}^3
 &=0.5443940332249657482536718308092\ldots,\\
 \mathbb EX_{1,0.005}^3
 &=0.5443321558595570316756456585706\ldots.
 \end{aligned}
\]
consistent with the proved limit
$2\sqrt6/9=0.5443310539518173551549520166013\ldots$.
These floating point residuals are independent representation checks,
not certified error intervals or replacements for the analytic proofs.

% END sections/01_joint_polygamma.tex

% BEGIN sections/02_gap_reconstruction.tex
% Parent-ready research fragment. Labels use eg:; no preamble is included.
% References to define in the final article: OrderedIncoming, GapSixIncoming,
% DLMFStirling. Sources inspected: current ProveIt incoming archives dated
% 2026-10-11 and 2026-10-12. The complete gap-six result is inherited.

\section{Effective reconstruction at every fixed gap}
\label{eg:section}

The incoming ordered-boundary report \cite{OrderedIncoming} constructs a Gamma integral basis
for the unequal-even-tail moment problem and proves eventual
nonsingularity for each fixed gap.  The subsequent gap-six report \cite{GapSix}
makes one of those eventual statements effective and completes that
ray by finite-field arithmetic.  We now give a fully explicit threshold
for \emph{every} fixed gap.  We then sharpen the estimates at gap eight
and certify all remaining indices.  In particular, the complete
gap-eight family, which the inspected gap-six report explicitly leaves
open, is resolved below.

The conclusions are statements about rational matrices on specified
formal labels.  They do not imply algebraic or rational independence
of the evaluated zeta values.  The contribution relative to the
inspected sources is the explicit general threshold and the completed
gap-eight ray; no historical-priority claim is made.

\subsection{The inherited rational obstruction}

Fix integers $m\ge1$ and $b>m$, and put
\[
 p=2b-2m,\qquad q=2b,\qquad x_n=n+\tfrac12,
 \qquad A_v(n)=\zeta(v,n+1).
\]
The meromorphic continuation of
\[
 D_{p,q}(s)=\sum_{n=0}^{\infty}A_p(n)A_q(n)x_n^{-s}
\]
is regular at every nonpositive even argument.  Indeed, the centered
Hurwitz expansion has only powers $x_n^{1-v-2h}$, so the possible
poles of the continued product sum are at the odd integers
$3-p-q-2h$.  Write $D_r=D_{p,q}(-2r)$.

The block coordinates from \cite{OrderedIncoming,GapSix} satisfy
\begin{align}
 D_r&=\sum_{j=0}^r\frac{\binom{2r}{2j}}{2j+1}
             B_{2r-2j}(1/2)C_j,\label{eg:forward}\\
 C_j&=\sum_{r=0}^j\binom{2j+1}{2r}2^{-2(j-r)}D_r.
 \label{eg:inverse}
\end{align}
For $j\ge b$ they have the late coordinates
\[
 C_j=\sum_{k=0}^{b-1}a_{jk}\zeta(2k),\qquad \zeta(0)=-\tfrac12.
\]
Put $d=2j+1$ and $e=d-p-q+2k+1$.  The coefficient is zero
if $e<0$; otherwise it is
\begin{equation}\label{eg:late-entry}
 a_{jk}=B_e\left\{
       \frac{\binom{d-q+1}{p-2k}}{d-q+1}
      +\frac{\binom{d-p+1}{q-2k}}{d-p+1}\right\}.
\end{equation}
Each summand is omitted unless its lower binomial index is between
$2$ and the upper index.  In particular the lower index zero is
omitted.  Define
\[
 M_{b,m}=(a_{b+r,k})_{0\le r,k<b}.
\]
The first $2b$ moment rows are expressed on the formal labels
\[
 C_0,\ldots,C_{b-1},\quad\zeta(0),\ldots,\zeta(2b-2).
\]
Their exact coefficient determinant is $\det M_{b,m}/(4b-1)!!$.
Thus nonsingularity of $M_{b,m}$ is exactly the optimal-prefix question.

We recall the fixed-gap reduction in a notation which keeps the
Bernoulli coefficients distinct from recurrence factors.  Set
\begin{align*}
 \beta_n&=\frac{2|B_{2n}|}{(2n)!},\qquad n\ge1,\\
 a_u&=0\quad(1\le u\le m),\qquad
 a_u=(2u)^{\underline{2m}}\quad(u>m),\\
 b_s&=(2s-1)^{\underline{2m}},\qquad
 d_s=2(a_s+b_s)\quad(s>m).
\end{align*}
Here the underline means falling factorial.  Define the rational
fundamental array by
\begin{align}
 R_{u,j}&=\delta_{uj}&&(1\le u,j\le m),\notag\\
 R_{s,j}&=\frac1{d_s}\sum_{u=1}^{s-1}
      \beta_{s-u}(b_s+a_u)R_{u,j}\quad&&(s>m).
 \label{eg:recurrence}
\end{align}
Finally let
\[
 F_{b,m}=(R_{b+i,j})_{1\le i,j\le m},\qquad f_{b,m}=\det F_{b,m}.
\]
The inherited complementary-minor identity is
\begin{equation}\label{eg:rational-bridge}
 \det M_{b,m}=
 \frac{(-1)^{b(b-1)/2}}{4^bb!}
       \left(\prod_{s=m+1}^{b+m}d_s\right) f_{b,m}.
\end{equation}
Every factor in the product is strictly positive.  Its normalization
includes the omitted-index convention in \eqref{eg:late-entry}.
This identity follows by coefficient comparison in the null equation
\[
 (R^{(2m)}-R^{(2m)}(0))\cot(z/2)
                         +(R\cot(z/2))^{(2m)}=0
\]
and the complementary-minor identity for its lower-triangular
coefficient operator; these inherited steps are reproduced in the
ordered-boundary and gap-six sources.

\subsection{Gamma realization and its signs}

Put
\begin{align*}
 W_m(x)&=\left|\Gamma\!\left(\tfrac14+\tfrac{ix}{2}\right)\right|^{4m},\\
 P_j(x)&=\sinh(\pi(m-j)x)\quad(1\le j<m),\qquad
 P_m(x)=\tanh(\pi x),\\
 \lambda_j&=j\pi,\qquad A_m=2^m(2\pi)^{2m},\qquad
 a_j^*=\begin{cases}A_m/2,&j<m,\\ A_m,&j=m.
                   \end{cases}
\end{align*}
The incoming Gamma-basis theorem \cite{OrderedIncoming} gives
\begin{equation}\label{eg:gamma-coefficients}
 t_{n,j}=\frac4{(2n)!}\operatorname{Im}
       \int_0^\infty(x+i/2)^{2n}P_j(x)W_m(x)\,dx
\end{equation}
and
\begin{equation}\label{eg:gamma-bridge}
 (t_{b+i,j})_{i,j=1}^m=F_{b,m}T_m,\qquad
 T_m=(t_{n,j})_{n,j=1}^m,\qquad
 \operatorname{sgn}\det T_m=(-1)^{m(m-1)/2}.
\end{equation}
All integrals converge absolutely.  The null equation follows from
the Gamma recurrence after two Fourier contours are shifted by
$\pm i$; the sinh modes cross no poles, and the tanh mode supplies
the two residues at $\pm i/2$ which give the subtracted derivative.
The sign in \eqref{eg:gamma-bridge} follows from continuous
Cauchy--Binet: its integrand is a positive weight times the product
of a Vandermonde in $x_r^2$ and a reversed Vandermonde in
$\cosh(\pi x_r)$.  These exact inputs, rather than numerical
normalizations of the Gamma integrals, will be used below.

\subsection{An explicit real-density estimate for arbitrary gaps}

\begin{lemma}\label{eg:density}
For $X=2m$, every $x\ge X$, and $1\le j\le m$,
\begin{equation}\label{eg:density-bound}
 \frac{P_j(x)W_m(x)}{a_j^*}
   =x^{-m}e^{-\lambda_jx}(1+\delta_{m,j}(x)),
 \qquad |\delta_{m,j}(x)|\le\frac{2m}{x^2}.
\end{equation}
On $0\le x\le X$ one has the elementary bound
\begin{equation}\label{eg:compact-density}
 0\le P_j(x)W_m(x)/a_j^*<2^{16m^2+12m}.
\end{equation}
\end{lemma}

\begin{proof}
The complex Stirling remainder \cite[\S5.11(ii)]{DLMF}, after retaining $1/(12z)$, gives
\[
 \log\Gamma(z)=(z-\tfrac12)\log z-z+\tfrac12\log(2\pi)
                        +\frac1{12z}+\rho(z),\qquad
 |\rho(z)|\le\frac1{90|z|^3}
\]
for $z=1/4+ix/2$.  Indeed the first omitted term has magnitude
$1/(360|z|^3)$ and the secant factor is at most four in this
closed right-half-plane sector.  Let $t=1/(2x)$ and
\[
 L_m(x)=\log\frac{W_m(x)}{A_mx^{-m}e^{-m\pi x}}.
\]
Taking $4m$ times the real part gives exactly
\[
 L_m(x)=-\frac m2\log(1+t^2)
      +m\left(\frac{\arctan t}{t}-1\right)
      +\frac{4m}{3}\frac{t^2}{1+t^2}+4m\operatorname{Re}\rho(z).
\]
Using $\log(1+t^2)\le t^2$ and
$0\le1-\arctan(t)/t\le t^2/3$, for $x\ge2$ we obtain
\begin{equation}\label{eg:log-density}
 |L_m(x)|\le\frac{13m}{24x^2}+\frac{16m}{45x^3}
          \le\frac{259m}{360x^2}<\frac{3m}{4x^2}.
\end{equation}
For $x\ge2m$ this is at most $3/(16m)<1$.  The elementary
inequality $|e^u-1|\le2|u|$ for real $|u|\le1$ now gives
an error at most $3m/(2x^2)$ for the normalized Gamma density.
The remaining normalized sinh or tanh multiplier lies in $[0,1]$
and differs from one by at most $2e^{-2\pi x}\le1/(9x^2)$.
Combining these estimates proves \eqref{eg:density-bound}.

Euler's integral gives $|\Gamma(1/4+ix/2)|\le\Gamma(1/4)<5$.
For $0\le x\le2m$, each multiplier is at most
$e^{2\pi m(m-1)}$.  Since $a_j^*>1$, $\pi<4$, $e<4$,
and $5<8$, the stated generous power-of-two bound follows.
\end{proof}

\subsection{An effective shifted Laplace lemma}

For $N>m+1$ and $\lambda>0$ define
\[
 S_N^{(r)}(\lambda)=
       \frac{\Gamma(N-m-r+1)}{\lambda^{N-m-r+1}},
       \qquad r=0,1,2.
\]
Let
\begin{equation}\label{eg:constants}
 C_m=\left\lceil4m+\frac{3m(m+1)}8+1\right\rceil,
 \qquad N_0(m)=4096(m+1)^2.
\end{equation}

\begin{lemma}\label{eg:laplace}
For every integer $N\ge N_0(m)$ and $1\le j\le m$,
\begin{equation}\label{eg:laplace-bound}
 \left|\frac1{a_j^*}\int_0^\infty(x+i/2)^NP_j(x)W_m(x)\,dx
   -e^{i\lambda_j/2}\left(S_N^{(0)}(\lambda_j)
                 +\frac{im}{2}S_N^{(1)}(\lambda_j)\right)\right|
       < C_m S_N^{(2)}(\lambda_j).
\end{equation}
\end{lemma}

\begin{proof}
Write $X=2m$.  On $[X,\infty)$ replace the normalized density
by $x^{-m}e^{-\lambda x}$.  Its absolute error is at most
\[
 2m\int_X^\infty |x+i/2|^Nx^{-m-2}e^{-\lambda x}\,dx.
\]
On $x\ge\sqrt N$ one has
$|x+i/2|^N\le x^N\exp(N/(8x^2))<2x^N$, so this part is
at most $4mS_N^{(2)}$.  On $X\le x\le\sqrt N$, the factor
$2mx^{-m-2}$ is at most one; hence the remaining part is at most
$\sqrt N(N+1)^{N/2}$.

For the model integral set $y=x+i/2$ and move the ray to
$[X,\infty)$.  Its integrand, after extracting $e^{i\lambda/2}$,
is $y^{N-m}(1-i/(2y))^{-m}e^{-\lambda y}$.  No pole is crossed.
On real $y\ge2m$, Taylor's theorem applied to $(1-u)^{-m}$ gives
\[
 \left|(1-i/(2y))^{-m}-1-\frac{im}{2y}\right|
       \le\frac{3m(m+1)}{8y^2}.
\]
For detail, $u=1/(2y)\le1/(4m)$, and the positive-coefficient
Taylor remainder is bounded by
$m(m+1)u^2(1-u)^{-m-2}/2$.  Also
\[
 (1-u)^{-m-2}\le
  \exp\!\left(\frac{(m+2)u}{1-u}\right)
  \le\exp\!\left(\frac{m+2}{4m-1}\right)<3.
\]
Thus the model remainder contributes at most
$3m(m+1)S_N^{(2)}/8$.

The short vertical segment at $X$, the two retained integrals
between $0$ and $X$, and the original compact density integral
are together bounded by $2^{20(m+1)^2}(3m)^N$.  Indeed the
compact density bound contributes at most
$X2^{16m^2+12m}(X+1)^N$; the other three terms are bounded
by $(2+m/2)(X+1)^N$.  The displayed larger constant covers their
sum for every $m\ge1$.

It remains to dominate
\[
 R_N=2^{20(m+1)^2}(3m)^N+\sqrt N(N+1)^{N/2}
\]
by $S_N^{(2)}$.  Put $k=N-m-2$.  Since $\lambda_j<4m$ and
$k!\ge(k/3)^k$,
\[
 S_N^{(2)}(\lambda_j)>
             \frac1{4m}\left(\frac{k}{12m}\right)^k.
\]
For $N\ge N_0(m)$, one has $k\ge3N/4$ and
$\sqrt{N+1}\le(5/4)\sqrt N$.  Consequently
\begin{align*}
 \frac{\sqrt N(N+1)^{N/2}}{S_N^{(2)}}
  &\le N^{m+4}\left(\frac{20m}{\sqrt N}\right)^N
    <N^{m+4}2^{-N},\\
 \frac{2^{20(m+1)^2}(3m)^N}{S_N^{(2)}}
  &\le2^{20(m+1)^2}N^{m+3}
                \left(\frac{48m^2}{N}\right)^N
    <2^{20(m+1)^2}N^{m+3}4^{-N}.
\end{align*}
The elementary inequality $N^{m+4}\le2^{N-2}$ holds throughout
this range.  At $N_0(m)$ it follows from
$\log_2N_0(m)\le12+2m$ and
$(m+4)(12+2m)<N_0(m)-2$; it persists by induction because
$N>2(m+4)$ implies $(1+1/N)^{m+4}<2$.
The two ratios are therefore less than $1/4$ each, the second
using also $N>20(m+1)^2$.  Thus $R_N<S_N^{(2)}/2$.
Adding all errors proves \eqref{eg:laplace-bound}.
\end{proof}

\subsection{A closed, computable threshold for every fixed gap}

Let $V_m=(j^{-2i})_{1\le i,j\le m}$ and define the positive rational
number
\begin{equation}\label{eg:delta}
 \delta_m=\frac{|\det V_m|}{\operatorname{perm}V_m}.
\end{equation}
Here $\operatorname{perm}V_m$ is the sum of its $m!$ unsigned
permutation products.  Thus $0<\delta_m\le1$, and its numerator
is explicitly the absolute reciprocal-square Vandermonde product
\[
 |\det V_m|=
 \left(\prod_{j=1}^m j^{-2}\right)
             \prod_{j<k}(j^{-2}-k^{-2}).
\]
Set
\begin{equation}\label{eg:B}
 E_m=8C_m+2(m-1),\qquad
 \boxed{\displaystyle
 B(m)=\max\left\{2048(m+1)^2,
       1+\left\lceil\frac{2mE_m}{\delta_m}\right\rceil\right\}.}
\end{equation}

\begin{theorem}[Effective reconstruction on every fixed-gap ray]
\label{eg:effective}
For every integer $m\ge1$ and every $b\ge B(m)$,
\begin{equation}\label{eg:general-sign}
 f_{b,m}\ne0,\qquad
 \operatorname{sgn}f_{b,m}=(-1)^{\lfloor m^2/4\rfloor}.
\end{equation}
Consequently the first $q=2b$ moments for
$(p,q)=(2b-2m,2b)$ form an optimal formal prefix, and
\[
 \operatorname{sgn}\det M_{b,m}
       =(-1)^{b(b-1)/2+\lfloor m^2/4\rfloor}.
\]
For any prescribed $m$, all possible exceptions lie in the explicitly
specified finite interval $m+1\le b<B(m)$.
\end{theorem}

\begin{proof}
Let $N_i=2b+2i$ for $1\le i\le m$.  At odd $j$, the leading
phase in \eqref{eg:laplace-bound} is $\pm i$, and at even $j$
it is $\pm1$.  Thus the odd-mode leading term is
\[
 \sigma_j\frac{4a_j^*}{N_i!}S_{N_i}^{(0)}(\lambda_j),
 \qquad \sigma_j=\sin(j\pi/2),
\]
with relative error bounded by
\[
 \frac{C_m\lambda_j^2}{(N_i-m)(N_i-m-1)}
       <\frac{16C_mm^2}{(N_i-m-1)^2}.
\]
At even $j$ the first nonzero imaginary term is
\[
 \sigma_j\frac{2ma_j^*}{N_i!}S_{N_i}^{(1)}(\lambda_j),
 \qquad \sigma_j=\cos(j\pi/2),
\]
with relative error less than $8C_m/(N_i-m-1)$.
Retaining this second phase term is essential.

Extract the positive row factors
\[
 \frac{4\Gamma(N_i-m+1)}{\Gamma(N_i+1)}\pi^{-2i}
\]
and the positive column factors $a_j^*\lambda_j^{-2b+m-1}$,
multiplying each even column factor additionally by
$m\lambda_j/(2(2b+2-m))$.  Also extract its sign $\sigma_j$.
The remaining array $E=(E_{ij})$ has the form
\[
 E_{ij}=j^{-2i}(1+\eta_{ij}).
\]
For an even column the extra row factor is
\[
 \rho_i=\frac{2b+2-m}{2b+2i-m},\qquad
 1-\rho_i\le\frac{2(m-1)}{2b+2-m}.
\]
Since $b\ge B(m)>m$ and $b\ge2m^2$, the preceding estimates give
\begin{equation}\label{eg:eta}
 |\eta_{ij}|<\frac{E_m}{b}<\frac{\delta_m}{2m}.
\end{equation}
For the odd columns use
$16C_mm^2/b^2\le8C_m/b$; for the even columns use
$|\rho_i(1+\epsilon)-1|\le1-\rho_i+|\epsilon|$.

Expanding the determinant into its signed permutation products,
\[
 |\det E-\det V_m|
 <\operatorname{perm}V_m
       \left\{\left(1+\frac{\delta_m}{2m}\right)^m-1\right\}
 <\operatorname{perm}V_m\,\delta_m=|\det V_m|.
\]
The last strict inequality holds for $0<\delta_m\le1$: the binomial
sum is bounded above by the corresponding infinite geometric sum
$\delta_m/(2-\delta_m)\le\delta_m$, with a strictly positive
omitted tail.  Hence $\det E$ has the sign of the reversed
Vandermonde, namely $(-1)^{m(m-1)/2}$.
The extracted column signs have product
$\prod_j\sigma_j=(-1)^{\lfloor m^2/4\rfloor}$.
Dividing by the initial determinant in \eqref{eg:gamma-bridge}
cancels the reversed-Vandermonde sign and proves
\eqref{eg:general-sign}.  Apply \eqref{eg:rational-bridge} to
obtain the assertion about the original moment matrix.
\end{proof}

The formula is deliberately conservative and does not claim a least
threshold.  Its first values are
\[
 \begin{array}{c|rrrrrr}
 m&1&2&3&4&5&6\\ \hline
 B(m)&8192&18432&32768&51200&635656&16146233.
 \end{array}
\]
The earlier complete gap-two, gap-four, and gap-six theorems are
stronger on their individual rays.  The advance here is an effective
formula valid for arbitrary prescribed $m$.  It does not assert
that every one of the finite remaining matrices is nonsingular;
the unrestricted all-pairs conjecture remains open.

\section{The complete gap-eight theorem}
\label{eg:gap-eight}

For $m=4$ sharper constants reduce the finite-completion threshold
from $51200$ to $8192$.  The next two estimates supply that reduction.

\begin{lemma}[Gap-eight density and phase estimates]
\label{eg:eight-laplace}
For $m=4$, $x\ge2$, and $1\le j\le4$,
\[
 \frac{P_j(x)W_4(x)}{a_j^*}
    =x^{-4}e^{-\lambda_jx}(1+\delta_j(x)),\qquad
 |\delta_j(x)|<\frac5{x^2}.
\]
For every integer $N\ge16384$,
\begin{equation}\label{eg:eight-phase}
 \left|\frac1{a_j^*}\int_0^\infty(x+i/2)^NP_j(x)W_4(x)\,dx
 -e^{i\lambda_j/2}
     \left(\frac{\Gamma(N-3)}{\lambda_j^{N-3}}
            +2i\frac{\Gamma(N-4)}{\lambda_j^{N-4}}\right)\right|
       <16\frac{\Gamma(N-5)}{\lambda_j^{N-5}}.
\end{equation}
\end{lemma}

\begin{proof}
Equation \eqref{eg:log-density} gives $|L_4(x)|<3/x^2\le3/4$.
The elementary Taylor-series bound $e^{3/4}<17/8$ implies
$|e^u-1|\le(3/2)|u|$ for real $|u|\le3/4$.  Adding the
normalized multiplier error $1/(9x^2)$ gives
\[
 |\delta_j(x)|<\left(\frac92+\frac19\right)x^{-2}<5x^{-2}.
\]
On $0\le x\le2$ the normalized density is less than $2^{80}$,
since $5^{16}<2^{38}$, $e^{6\pi}<3^{24}<2^{39}$, and $a_j^*>1$.

Repeat the proof of Lemma~\ref{eg:laplace} with $X=2$.  The
density error above $\sqrt N$ contributes less than
$10\Gamma(N-5)/\lambda_j^{N-5}$.  The model Taylor remainder
has the sharper bound
\[
 \left|(1-i/(2y))^{-4}-1-\frac{2i}{y}\right|<\frac5{y^2}
       \qquad(y\ge2).
\]
Indeed the positive series remainder divided by $u^2$, where
$u=1/(2y)\le1/4$, is increasing and equals $1504/81$ at $u=1/4$;
its contribution is therefore $376/(81y^2)<5/y^2$.

All compact and lower-range contributions are at most
\[
 R_N=2^{82}3^N+\sqrt N(N+1)^{N/2}.
\]
Now $\lambda_j\le4\pi<13$ and, with $k=N-6$,
\[
 \frac{\Gamma(N-5)}{\lambda_j^{N-5}}
       >\frac1{13}\left(\frac{k}{39}\right)^k.
\]
For every $N\ge16384$,
\[
 \frac{39\sqrt{N+1}}{N-6}<\frac12,\qquad
 \frac{117}{N-6}<\frac1{128},\qquad N^7\le2^{N-2}.
\]
The first inequality follows by squaring, checking the resulting
quadratic at $16384$, and noting its positive derivative thereafter;
the second is immediate; the third holds at $16384$ and persists
by the same induction as in Lemma~\ref{eg:laplace}.  Therefore
the two normalized parts of $R_N$ are bounded, respectively, by
$2^{82}N^6 2^{-7N}<1/4$ and $N^7 2^{-N}\le1/4$.
The total error is consequently less than $(10+5+1/2)$ times the
displayed Gamma factor, proving \eqref{eg:eight-phase}.
\end{proof}

\begin{proposition}[Certified positive tail at gap eight]
\label{eg:eight-positive}
For every integer $b\ge8192$, $f_{b,4}>0$.
\end{proposition}

\begin{proof}
In \eqref{eg:eight-phase}, an odd column has relative error less
than $2704/(N-5)^2$, and an even column has relative error less
than $104/(N-5)$; these use only $\lambda_j<13$.
Perform the row and column normalizations in Theorem~\ref{eg:effective}.
The even-column row factors are now
$\rho_i=(2b-2)/(2b+2i-4)$, with
$1-\rho_i\le6/(2b-2)$.
Hence the normalized entries satisfy
\[
 E_{ij}=j^{-2i}(1+\eta_{ij}),\qquad
 |\eta_{ij}|<\frac1{128}\qquad(b\ge8192).
\]
For the even columns use
$6/(2b-2)+104/(2b-3)<110/(2b-3)<1/128$;
the odd-column inequality is weaker still.

The four-by-four limiting matrix has exact values
\[
 \det V_4=\frac{175}{127401984},\qquad
 \operatorname{perm}V_4=\frac{157183}{4586471424},\qquad
 \frac{\det V_4}{\operatorname{perm}V_4}=\frac{6300}{157183}.
\]
Moreover
\[
 \det V_4-\operatorname{perm}V_4
        \left\{(129/128)^4-1\right\}
       =\frac{357064449025}{1231171548132409344}>0.
\]
Thus $\det E>0$.  The extracted signs are $(+,-,-,+)$, whose
product is positive, and $\det T_4>0$.  Equation
\eqref{eg:gamma-bridge} proves $f_{b,4}>0$.
\end{proof}

\subsection{An exact finite-field completion}

The cotangent coefficients admit the elementary recurrence
\begin{equation}\label{eg:beta-recurrence}
 \beta_1=\frac16,\qquad
 \beta_n=\frac1{2(2n+1)}\sum_{k=1}^{n-1}\beta_k\beta_{n-k}
                  \quad(n\ge2).
\end{equation}
It follows by substituting the Taylor expansion of
$(z/2)\cot(z/2)$ into its Riccati equation.
Equations \eqref{eg:beta-recurrence} and \eqref{eg:recurrence}
can be computed exactly modulo the prime
$\ell=2147483647=2^{31}-1$.

For $m=4$ and $5\le b\le8191$ the largest fundamental row
required is $8195$, so only $\beta_1,\ldots,\beta_{8194}$ are needed.
Every denominator is invertible modulo $\ell$.  The denominators
$2(2n+1)$ have their prime factors below $\ell$, and
\begin{equation}\label{eg:diagonal-factorization}
 d_s=2\frac{(2s)!}{(2s-8)!}\frac{2s-4}{s}\qquad(s\ge5)
\end{equation}
has the same property throughout the stated finite range.  Thus
induction identifies every computed residue with the reduction of
the exact rational entry.

\begin{proposition}[Finite gap-eight certificate]
\label{eg:eight-finite}
For every integer $5\le b\le8191$, the rational determinant
$f_{b,4}$ has a nonzero reduction modulo $2147483647$.
\end{proposition}

\begin{proof}[Exact computational proof]
The accompanying \texttt{verify\_gap\_eight.cpp} uses integer
arithmetic only.  It verifies primality by trial division through
$\lfloor\sqrt\ell\rfloor$, generates \eqref{eg:beta-recurrence},
generates the four fundamental columns from \eqref{eg:recurrence},
and evaluates each four-by-four determinant by finite-field Gaussian
elimination.  All $8187$ determinant residues are nonzero; the full
indexed list is supplied in \texttt{gap\_eight\_certificate.json}.
As an independent exact check of the cotangent normalization, the
program verifies all $8194$ coefficients by multiplying
$1-\tfrac12\sum_{n\ge1}\beta_nt^n$ by the directly generated
series of $\sin(\sqrt t/2)/(\sqrt t/2)$ and comparing with the
direct series of $\cos(\sqrt t/2)$.
The separate SymPy verifier builds rational Bernoulli numbers from
their standard definition and checks the original matrix
\eqref{eg:late-entry}, the full bridge \eqref{eg:rational-bridge},
and the finite-field residues independently for $5\le b\le16$.
These latter checks guard against a sign, factorial, or omitted-index
error in passing from the source matrix to the recurrence.
\end{proof}

\begin{theorem}[Every gap-eight pair has the optimal prefix]
\label{eg:eight-main}
For every integer $b\ge5$, the matrix $M_{b,4}$ is nonsingular.
For $p=2b-8$ and $q=2b$, consequently,
\begin{align*}
 \operatorname{span}_{\mathbb Q}\{D_0,D_1,\ldots\}
 &=\operatorname{span}_{\mathbb Q}\{D_0,\ldots,D_{q-1}\}\\
 &=\operatorname{span}_{\mathbb Q}
       \{C_0,\ldots,C_{b-1},\zeta(0),\ldots,\zeta(2b-2)\}.
\end{align*}
The first $q$ moment rows have formal rank $q$ and coefficient
determinant $\det M_{b,4}/(4b-1)!!$.  This is optimal for those
$q$ formal coordinates.
\end{theorem}

\begin{proof}
Proposition~\ref{eg:eight-positive} covers $b\ge8192$, and
Proposition~\ref{eg:eight-finite} covers $5\le b\le8191$.
Equation \eqref{eg:rational-bridge} proves nonsingularity of the
source matrix.  The triangular transforms
\eqref{eg:forward}--\eqref{eg:inverse} give the stated generation
and formal-rank conclusions.
\end{proof}

The certificate proves nonvanishing at the finite indices; it does
not determine their real signs.  Direct rational checks give
$f_{b,4}>0$ for $5\le b\le16$, and the analytic theorem gives
positivity for $b\ge8192$.  We make no positivity assertion for
the intervening finite interval on the basis of the residue test.

\subsection{An all-order family of absolutely convergent identities}

The completed ray produces ordinary convergent evaluations of
polygamma-tail products at every even pair separated by eight.
Let $(\lambda_b,\ldots,\lambda_{2b-1})$ be row zero of
$M_{b,4}^{-1}$ and define
\begin{align*}
 K_b(y)&=\sum_{j=b}^{2b-1}\lambda_jy^{2j+1},\\
 P_b(x)&=K_b(x+\tfrac12)-K_b(x-\tfrac12)
             =\sum_{r=0}^{q-1}w_rx^{2r}.
\end{align*}
For $0\le h\le4$ put
\[
 c_{v,0}=\frac1{v-1},\qquad
 c_{v,h}=\frac{B_{2h}(1/2)}{(2h)!}(v)_{2h-1}\quad(h\ge1),
 \qquad u_h=\sum_{e=0}^h c_{p,e}c_{q,h-e},
\]
and set
\[
 S_b(x)=\sum_{h=0}^4
       \left(\sum_{e=0}^h u_e w_{q-1-h+e}\right)x^{8-2h}.
\]
Every coefficient here is rational and specified by finite algebra.

\begin{corollary}[Rational polygamma-product evaluations]
\label{eg:eight-identities}
For every $b\ge5$, $p=2b-8$, $q=2b$, and $x_n=n+1/2$,
\begin{equation}\label{eg:eight-sum}
 \boxed{\displaystyle
 \sum_{n=0}^{\infty}\left\{
 \frac{P_b(x_n)}{(p-1)!(q-1)!}
       \psi^{(p-1)}(n+1)\psi^{(q-1)}(n+1)-S_b(x_n)
                         \right\}=-\frac12.}
\end{equation}
The sum converges absolutely.  The polynomial $P_b$ has exact degree
$2q-2$ and leading coefficient
\begin{equation}\label{eg:leading}
 w_{q-1}=\frac{4b(4b-1)f_{b-1,4}}{d_{b+4}f_{b,4}}\ne0,
\end{equation}
where at the initial index one uses
\[
 f_{4,4}=\frac1{13105577681191934125041254400000000}.
\]
\end{corollary}

\begin{proof}
By construction, $\sum_{j=b}^{2b-1}\lambda_j C_j=\zeta(0)=-1/2$.
The centered difference defining $P_b$ has coefficients
\[
 w_r=\sum_{j=\max\{b,r\}}^{2b-1}
       \lambda_j\binom{2j+1}{2r}2^{-2(j-r)}.
\]
Thus \eqref{eg:inverse} gives $\sum_rw_rD_r=-1/2$.
The first five terms of the centered tail product are
\[
 A_p(n)A_q(n)=\sum_{h=0}^4u_hx_n^{2-p-q-2h}
                      +O(x_n^{-p-q-8}).
\]
Multiplication by $P_b$ leaves precisely $S_b$ as its nonnegative
powers, with the remainder $O(x_n^{-2})$.  The analytically
continued sums of $1,x_n^2,\ldots,x_n^8$ vanish because
$\zeta(-2h,1/2)=0$.  Therefore the ordinary absolutely convergent
sum in \eqref{eg:eight-sum} equals that moment combination.
Finally, for even $v$,
$A_v(n)=\psi^{(v-1)}(n+1)/(v-1)!$.

For $b\ge6$, deleting the last row and first column of $M_{b,4}$
and reindexing $j,k,p,q$ downward by $1,1,2,2$, respectively,
gives exactly $M_{b-1,4}$.  Cramer's rule and
\eqref{eg:rational-bridge} give
$\lambda_{2b-1}=4b f_{b-1,4}/(d_{b+4}f_{b,4})$.
The centered difference multiplies the top coefficient by $4b-1$,
proving \eqref{eg:leading}.  At $b=5$ the same relation is
verified directly with the displayed rational auxiliary determinant.
Theorem~\ref{eg:eight-main} proves all other required nonvanishing.
\end{proof}


\subsection{A concrete pair beyond the previous small-order coverage}

For $(p,q)=(14,22)$, namely $b=11$, the construction gives an even
polynomial $P_{11}$ of degree $42$ and a subtraction polynomial
$S_{11}$ of degree $8$. The leading coefficient is
\[
 [x^{42}]P_{11}(x)=
 \frac{452253373522332335630380}{10060146014319869101877031}.
\]
All $22$ coefficients of $P_{11}$ and all five coefficients of
$S_{11}$ are supplied as exact rationals in
\texttt{gap\_eight\_14\_22\_example.json}.
The short generator \texttt{create\_gap\_eight\_example.py} checks
the inverse row against the original Bernoulli matrix and verifies
the entire subtraction polynomial symbolically. Thus the example
specifies a directly usable ordinary sum, without printing several
pages of large rational coefficients. The finite definitions above
are the full exact formula for every other pair as well.

% END sections/02_gap_reconstruction.tex

% BEGIN sections/03_multiscale_collisions.tex
\section{Stieltjes collisions at different powers of the small parameter}
\label{mc:section}

The unequal-rate collision theorem in the latest gap-six report
\cite[Section~4 and research question~10]{GapSix} assumes
$a=\lambda\epsilon$ and $b=\mu\epsilon$ with fixed positive rates.
Its proof explicitly leaves different powers of $\epsilon$ open.
Here we resolve that question for every pair of positive monomial
scales and every fixed pair of Stieltjes indices and argument
derivatives. The additional constants occur at explicitly classified
rational ratios of the exponents. In particular, the product of the
two singular parts can now contribute a constant.

The local subtraction argument is a classical asymptotic method;
related generalized Stieltjes-transform methods appear in
\cite{TicaGalapon}. The contribution here is the complete resonance
classification and its evaluated contact formula for this family,
rather than a claim to invent finite-part asymptotics.

\subsection{A compact identity before the general theorem}

Write $\gamma$ for Euler's constant, and put
\begin{equation}\label{mc:D}
 D(a,b)=\frac{\log(b/a)}{b-a}
       =\int_0^\infty\frac{dx}{(x+a)(x+b)},\qquad a,b>0.
\end{equation}
The value at $a=b$ is $1/a$. A derivative in the first argument is
written $D_a$; explicitly,
\begin{equation}\label{mc:Da}
 D_a(a,b)=-\frac1{a(b-a)}+
                    \frac{\log(b/a)}{(b-a)^2}.
\end{equation}

\begin{theorem}[A two-scale ordered digamma identity]\label{mc:ordered-limit}
Fix $\lambda,\mu,\alpha,\beta>0$, and set
$a=\lambda\epsilon^\alpha$, $b=\mu\epsilon^\beta$.
Then, as $\epsilon\downarrow0$,
\begin{equation}\label{mc:ordered-uniform}
 \boxed{\displaystyle
 \int_0^1\psi'(x+a)\psi(x+b)\,dx
 =D_a(a,b)+\frac{\psi(1+b)}a
                   +\zeta(2)\log\frac ba+o(1).}
\end{equation}
Every term on the left is an ordinary convergent integral. The
subtraction on the right is exact to a remainder tending to zero
along every indicated monomial path.
\end{theorem}

\begin{proof}
Let $h(x)=\psi(1+x)$, so that $\psi(x)=-1/x+h(x)$.
Expand the product and integrate the term
$h(x+b)/(x+a)^2$ by parts. The difference between the two sides of
\eqref{mc:ordered-uniform}, without $o(1)$, is
\begin{align}\label{mc:ordered-remainder}
 &-D_a(1+a,1+b)-\frac{h(1+b)}{1+a}
     +\zeta(2)\log\frac{1+a}{1+b}\notag\\
 &\quad+\int_0^1\left\{
       \frac{h'(x+b)-\zeta(2)}{x+a}
      -\frac{h'(x+a)-\zeta(2)}{x+b}\right\}dx
      +\int_0^1h'(x+a)h(x+b)\,dx.
\end{align}
In the middle integral replace $h'(x+b)$ and $h'(x+a)$ by
$h'(x)$. The two errors are
$O(b(1+|\log a|)+a(1+|\log b|))$.
Since $h'(x)-\zeta(2)=O(x)$, the remaining difference tends to
zero by dominated convergence. The first, second, and last terms
in \eqref{mc:ordered-remainder} tend respectively to
\[
 \frac12,\qquad \gamma-1,\qquad
 \frac12\bigl(\psi(2)^2-\psi(1)^2\bigr)=\frac12-\gamma.
\]
They cancel. A positive power of $\epsilon$ times any fixed power
of $|\log\epsilon|$ tends to zero, proving the statement.
\end{proof}

For a function with a power--logarithm expansion, let
$\operatorname{CT}_{\epsilon}$ denote the coefficient of
$\epsilon^0(\log\epsilon)^0$. It does not mean the ordinary limit
of a divergent integral. Let $\mathbf1_{\mathcal E}$ denote the
indicator of the stated condition.

\begin{corollary}[An explicit resonance law]\label{mc:ordered-CT}
Suppose $q=\alpha/\beta>1$, and set $j=2/(q-1)$.
For $I(\epsilon)$ equal to the integral in
\eqref{mc:ordered-uniform},
\begin{align}\label{mc:ordered-constant}
 \operatorname{CT}_{\epsilon}I
 ={}&\zeta(2)\log\frac\mu\lambda
  +\mathbf1_{q\in\mathbb Z_{\ge2}}
       (-1)^{q+1}\frac{\mu^q}{\lambda}\zeta(q+1)\notag\\
 &+\mathbf1_{j\in\mathbb Z_{\ge1}}
       \frac{\lambda^j}{\mu^{j+2}}
                  \left((j+1)\log\frac\mu\lambda-1\right).
\end{align}
A term whose indicator is zero is omitted, including its
integer-indexed notation. In particular, at $\lambda=\mu=1$,
\begin{equation}\label{mc:ordered-examples}
 \operatorname{CT}_{\epsilon}I=
 \begin{cases}
  -1-\zeta(3),&q=2,\\
  -1+\zeta(4),&q=3,\\
  (-1)^{q+1}\zeta(q+1),&q\in\mathbb Z,\ q\ge4,\\
  -1,&q=1+2/j,\ j\ge3,\\
  0,&q\notin\mathbb Z_{\ge2}\ \text{and }2/(q-1)\notin\mathbb Z_{\ge1}.
 \end{cases}
\end{equation}
\end{corollary}

\begin{proof}
Use the convergent expansion
$\psi(1+b)=-\gamma+\sum_{\ell\ge1}(-1)^{\ell+1}
\zeta(\ell+1)b^\ell$ for $|b|<1$.
In $\psi(1+b)/a$, exponent zero requires $\ell\beta=\alpha$.
If $\rho=a/b$, the other singular term has the convergent expansion
\begin{equation}\label{mc:Da-series}
 D_a(a,b)=-\frac1{ab}
 +\frac1{b^2}\sum_{k\ge0}\rho^k
                   \left((k+1)\log\frac ba-1\right).
\end{equation}
Here exponent zero requires $k(\alpha-\beta)=2\beta$.
Extracting the constant in $\log\epsilon$ proves the result.
\end{proof}

These are new exact values of specified asymptotic coefficients.
They also give ordinary renormalized limits by subtracting the
finitely many divergent powers and logarithms in
\eqref{mc:Da-series} and the displayed digamma series. For example,
if $L=\log\epsilon$, the first two cases are
\begin{align}\label{mc:q2limit}
 \lim_{\epsilon\downarrow0}\bigg[
  \int_0^1\psi'(x+\epsilon^2)\psi(x+\epsilon)dx
   +\epsilon^{-3}+(1+\gamma+L)\epsilon^{-2}\notag\\
 \hspace{28mm}
   +(1+2L-\zeta(2))\epsilon^{-1}
   +(3+\zeta(2))L\bigg]
 &=-1-\zeta(3),\\[3pt]
 \lim_{\epsilon\downarrow0}\bigg[
  \int_0^1\psi'(x+\epsilon^3)\psi(x+\epsilon)dx
   +\epsilon^{-4}+\gamma\epsilon^{-3}
   +(1+2L-\zeta(2))\epsilon^{-2}\notag\\
 \hspace{28mm}
   +\zeta(3)\epsilon^{-1}
   +(4+2\zeta(2))L\bigg]
 &=-1+\zeta(4).
 \label{mc:q3limit}
\end{align}
The convergence is established analytically by the theorem;
the numerical diagnostics in the package test the signs and
the independent expansions.

\subsection{Definitions for arbitrary Stieltjes indices}

We fix the standard convention
\begin{equation}\label{mc:stieltjes}
 Z(u,x)=\zeta(1+u,x)-\frac1u
     =\sum_{m\ge0}\frac{(-1)^m}{m!}\gamma_m(x)u^m,
 \qquad \gamma_0(x)=-\psi(x).
\end{equation}
Argument derivatives are indicated by a parenthesized superscript.
Put
\begin{equation}\label{mc:notation}
 \begin{gathered}
 P_r(u)=(-1)^r(1+u)_r,\qquad
 A_r(u)=\frac{(1+u)_r}{r!},\\
 Z_r(u,x)=\partial_x^r Z(u,x),\qquad
 g_h(u)=\left.\partial_x^h Z(u,x)\right|_{x=1}.
 \end{gathered}
\end{equation}
These are holomorphic near $u=0$. The Hurwitz recurrence
\cite[\S25.11]{DLMF} gives the exact splitting
\begin{equation}\label{mc:split}
       Z_r(u,x)=P_r(u)x^{-r-1-u}+Z_r(u,1+x).
\end{equation}
In particular $g_0(u)=\zeta(1+u)-1/u$ and, for $h\ge1$,
\begin{equation}\label{mc:g-explicit}
        g_h(u)=(-1)^h(1+u)_h\zeta(1+u+h).
\end{equation}
Thus every derivative of $g_h$ is a finite combination of
ordinary zeta derivatives and generalized harmonic numbers;
for $h=0$ its coefficients are the ordinary Stieltjes constants.

The reference finite part is
\begin{equation}\label{mc:FP}
 F_{m,n}^{r,t}=\operatorname{CT}_{\delta\downarrow0}
  \int_\delta^1\gamma_m^{(r)}(x)\gamma_n^{(t)}(x)\,dx.
\end{equation}
It uses the coordinate $\delta$ exactly as displayed. An explicit
holomorphic germ $\mathcal F_{r,t}(u,v)$ for these coefficients
can be defined entirely by an ordinary convergent integral.
With $R=r+t+1$, set
\begin{align}\label{mc:FP-remainder}
 \mathcal R(u,v;x)={}&Z_r(u,x)Z_t(v,x)
       -P_r(u)P_t(v)x^{-R-1-u-v}\notag\\
 &-P_r(u)\sum_{k=0}^r\frac{g_{t+k}(v)}{k!}x^{k-r-1-u}
  -P_t(v)\sum_{k=0}^t\frac{g_{r+k}(u)}{k!}x^{k-t-1-v},\\
 \mathcal F_{r,t}(u,v)={}&\int_0^1\mathcal R(u,v;x)dx
  -\frac{P_r(u)P_t(v)}{R+u+v}\notag\\
 &-P_r(u)\sum_{k=0}^{r-1}\frac{g_{t+k}(v)}{k!(r-k+u)}
  -P_t(v)\sum_{k=0}^{t-1}\frac{g_{r+k}(u)}{k!(t-k+v)}.
 \label{mc:FP-germ}
\end{align}
Empty sums vanish. Then
$F_{m,n}^{r,t}=(-1)^{m+n}\partial_u^m\partial_v^n
\mathcal F_{r,t}(0,0)$.
The critical powers $x^{-1}\log^h x$ contribute zero to the
constant on $(0,1)$; this explains their omission from the two
rational sums in \eqref{mc:FP-germ}. This construction agrees with
the inherited convention \cite{GapSix,LatticeIncoming}.

For example, those sources give
\begin{equation}\label{mc:FP-values}
 \begin{gathered}
 F_{0,0}^{0,0}=2\gamma_1-2\zeta(2),\quad
 F_{0,0}^{1,0}=\zeta(2),\quad
 F_{0,0}^{1,1}=2\zeta(3)+4\zeta'(3),\\
 F_{1,0}^{0,0}=\tfrac32\gamma_2+2\gamma\zeta(2)-\zeta(3).
 \end{gathered}
\end{equation}
These baseline evaluations are credited to the previous work.
The new result evaluates every change caused by the different scales.

\subsection{The complete constant formula}

Assume henceforth that $\alpha>\beta>0$ and define
\begin{equation}\label{mc:scales}
 a=\lambda\epsilon^\alpha,\quad b=\mu\epsilon^\beta,
 \quad d=\alpha-\beta,\quad \tau=\lambda/\mu,
 \quad R=r+t+1.
\end{equation}
The opposite ordering follows by exchanging the factors.
Introduce the finite set
\begin{equation}\label{mc:resonance-set}
 \mathscr L_r(\alpha,\beta)
 =\{(k,\ell):0\le k<r,\ \ell\in\mathbb Z_{\ge1},\
                      \beta\ell=\alpha(r-k)\}.
\end{equation}
It has at most $r$ members, including at rational ratios with
large numerator or denominator. Define the following removable germ:
\begin{align}\label{mc:mixed-contact}
 \mathcal C_{r,t}(u,v;\lambda,\mu;\alpha,\beta)
 ={}&(-1)^r g_{r+t}(v)
                  \frac{\lambda^{-u}-A_r(u)}u\notag\\
 &+(-1)^r\sum_{(k,\ell)\in\mathscr L_r(\alpha,\beta)}
   \frac{\mu^\ell\lambda^{k-r-u}}{\ell!}
                 (1+u)_{r-k-1}\,g_{t+k+\ell}(v).
\end{align}
The first quotient is removable since $A_r(0)=1$.

There is one possible singular--singular contact. If
$j=\beta R/d$ is a positive integer, set
\begin{align}\label{mc:singular-contact}
 \mathcal H_j(u,v)={}&
 \frac{(-1)^t\mu^{-R-v}\tau^j}{u\Gamma(1+v)}\notag\\
 &\quad\times\left\{
 \frac{\lambda^{-u}\Gamma(R+j+v)}{(1-u)_j}
 -\frac{\mu^{-u}\Gamma(1-u)\Gamma(R+j+u+v)}{j!}
                 \right\}.
\end{align}
The numerator vanishes identically at $u=0$, so this is a
holomorphic two-variable germ. Define $\mathcal H_*=\mathcal H_j$
when $\beta R/d=j\in\mathbb Z_{\ge1}$, and $\mathcal H_*=0$
otherwise. The quotient must be removed before taking derivatives.

\begin{theorem}[All-index monomial collision constants]\label{mc:main}
For every $m,n,r,t\in\mathbb Z_{\ge0}$ and all positive parameters
in \eqref{mc:scales}, the integral
\[
 I_{m,n}^{r,t}(\epsilon)
   =\int_0^1\gamma_m^{(r)}(x+a)
                   \gamma_n^{(t)}(x+b)\,dx
\]
has a power--logarithm asymptotic expansion through its constant.
Its exact constant is
\begin{align}\label{mc:main-constant}
 \boxed{\begin{aligned}
 \operatorname{CT}_{\epsilon}I_{m,n}^{r,t}
  =F_{m,n}^{r,t}+(-1)^{m+n}\partial_u^m\partial_v^n
    \big[&\mathcal C_{r,t}(u,v;\lambda,\mu;\alpha,\beta)\\
       +&\mathcal C_{t,r}(v,u;\mu,\lambda;\beta,\alpha)
         +\mathcal H_*(u,v)\big]_{u=v=0}.
 \end{aligned}}
\end{align}
Every correction in this formula is a finite expression in
logarithms, ordinary zeta values and derivatives, Stieltjes constants
at one, and generalized harmonic numbers. There is no new
unevaluated spectral constant in the correction.
\end{theorem}

There are precisely three possible mechanisms for extra constants:
\begin{equation}\label{mc:three-mechanisms}
 \beta\ell=\alpha(r-k),\qquad
 \alpha\ell=\beta(t-k),\qquad
 j(\alpha-\beta)=\beta(r+t+1).
\end{equation}
The two first equations have the finite ranges specified above;
the last has at most one admissible integer $j$. If $\alpha/\beta$
is irrational, all three resonance contributions vanish. Only the
first, critical quotient in each $\mathcal C$ remains. These are
selection rules for an asymptotic coefficient, not statements of
discontinuity of the original positive-shift integral.

\subsection{A two-shift expansion with analytic coefficients retained}

We prove the theorem in two steps. Keeping the analytic coefficients
at their actual small shifts is the useful replacement for a
fixed-ratio Taylor truncation. Define
\begin{equation}\label{mc:kernel}
 \mathcal B_{r,t}(u,v;a,b)
  =P_r(u)P_t(v)\int_0^\infty
             (x+a)^{-r-1-u}(x+b)^{-t-1-v}\,dx.
\end{equation}
This is an ordinary absolutely convergent integral for $u,v$ near
zero. It also has the elementary incomplete-beta representation
used below.

\begin{lemma}[Uniform germ along monomial paths]\label{mc:uniform-germ}
For fixed exponents and rates as in \eqref{mc:scales}, locally
uniformly on some complex neighborhood of $(u,v)=(0,0)$,
\begin{align}\label{mc:uniform-expansion}
 &\int_0^1 Z_r(u,x+a)Z_t(v,x+b)dx
    =\mathcal F_{r,t}(u,v)+\mathcal B_{r,t}(u,v;a,b)\notag\\
 &\quad+(-1)^r\sum_{k=0}^{r-1}
     (1+u)_{r-k-1}a^{k-r-u}Z_{t+k}(v,1+b)
   +(-1)^r g_{r+t}(v)\frac{a^{-u}-A_r(u)}u\notag\\
 &\quad+(-1)^t\sum_{k=0}^{t-1}
     (1+v)_{t-k-1}b^{k-t-v}Z_{r+k}(u,1+a)
   +(-1)^t g_{r+t}(u)\frac{b^{-v}-A_t(v)}v+o(1).
\end{align}
All fixed spectral derivatives of the remainder tend to zero.
\end{lemma}

\begin{proof}
Use \eqref{mc:split}. The singular--singular product on $(0,1)$
equals \eqref{mc:kernel} minus its tail on $(1,\infty)$; the latter
tends uniformly to $P_r(u)P_t(v)/(R+u+v)$.
For the first mixed product, use Taylor's theorem in $x$ at the
point $1+b$:
\[
 Z_t(v,1+b+x)=\sum_{k=0}^r
          \frac{Z_{t+k}(v,1+b)}{k!}x^k+O(x^{r+1}).
\]
The remainder has a bound uniform for small $b$ and for $v$ in a
fixed disk. After multiplication by the singular factor it is
bounded by $Cx^{-\eta}$, with $\eta<1$, independent of the
relative sizes of $a,b$. Dominated convergence therefore applies.
For $k<r$ the beta integral gives
\begin{align*}
 P_r(u)\int_0^1x^k(x+a)^{-r-1-u}dx
 ={}&(-1)^rk!(1+u)_{r-k-1}a^{k-r-u}\\
 &-\frac{P_r(u)}{r-k+u}+o(1).
\end{align*}
For $k=r$ its removable version is
\[
 P_r(u)\int_0^1x^r(x+a)^{-r-1-u}dx
   =(-1)^rr!\frac{a^{-u}-A_r(u)}u+o(1).
\]
One can first take $\Re u>0$, subtract the tail, and continue;
the divided difference is uniformly holomorphic near zero.
Replacing the coefficient of the critical term at $1+b$ by its
value at one costs $O(b\,a^{-\eta}(1+|\log a|))$ on a small disk.
Choose $\eta>0$ so small that $\beta-\alpha\eta>0$; the error
then tends to zero. For the exchanged term also impose
$\alpha-\beta\eta>0$. The upper endpoints and the remaining
integrals give exactly \eqref{mc:FP-germ}. This proves local
uniformity. Cauchy's integral formula on a slightly larger disk
justifies every prescribed spectral derivative.
\end{proof}

The Taylor expansion
\[
 Z_{t+k}(v,1+b)=\sum_{\ell\ge0}
                       g_{t+k+\ell}(v)\frac{b^\ell}{\ell!}
\]
now shows that a noncritical mixed term has base exponent
$\alpha(k-r)+\beta\ell$. Taking it far enough to have a positive
remainder exponent is always a finite task. Exponent zero gives
$\mathscr L_r(\alpha,\beta)$; the critical quotient contributes
its constant in $\log\epsilon$ by replacing $a^{-u}$ with
$\lambda^{-u}$. This proves exactly the two $\mathcal C$ terms.

\subsection{Pairing the singular branches before differentiation}

Put $A=r+1+u$, $B=t+1+v$, $C=A+B-1=R+u+v$ and
$\rho=a/b$. For $0<a<b$, the substitution
$y=(x+a)/(x+b)$ gives
\begin{equation}\label{mc:incomplete-beta}
 \frac{\mathcal B_{r,t}}{P_r(u)P_t(v)}
  =(b-a)^{-C}\left\{
     \frac{\Gamma(1-A)\Gamma(C)}{\Gamma(B)}
      -\int_0^\rho y^{-A}(1-y)^{C-1}\,dy\right\}.
\end{equation}
At generic nonintegral $A$ the last integral is understood by its
analytic continuation; the original integral in
\eqref{mc:kernel} fixes the branch and the removable values.
Expansion of this incomplete beta function, or the usual
hypergeometric connection formula \cite[\S15.8]{DLMF}, gives
two normally convergent series for $0<\rho<1$:
\begin{align}\label{mc:two-branches}
 \frac{\mathcal B_{r,t}}{P_r(u)P_t(v)}
 ={}&b^{-C}\frac{\Gamma(1-A)}{\Gamma(B)}
             \sum_{j\ge0}\frac{\Gamma(C+j)}{j!}\rho^j\notag\\
 &+a^{1-A}b^{-B}
      \sum_{k\ge0}(-1)^k(B)_k
                         \frac{\Gamma(A-1-k)}{\Gamma(A)}\rho^k.
\end{align}
The individual series have poles at $u=0$. Their sum is
holomorphic, and the poles must be paired at equal base exponents.

The $k<r$ terms of the second branch are holomorphic and have
base exponents
$-\alpha r-\beta(t+1)+k(\alpha-\beta)<0$.
For each $j\ge0$, pair the first-branch term indexed by $j$ with
the second-branch term indexed by $k=r+j$. Their common base
exponent is
\begin{equation}\label{mc:paired-exponent}
                    E_j=-\beta R+j(\alpha-\beta).
\end{equation}
After the Gamma reflection and recurrence identities, the paired
term is $\epsilon^{E_j}\mathcal H_j(u,v;L)$, where $L=\log\epsilon$
and
\begin{align}\label{mc:H-log}
 \mathcal H_j(u,v;L)={}&
 \frac{(-1)^t\mu^{-R-v}\tau^j\epsilon^{-\beta v}}
                   {u\Gamma(1+v)}\notag\\
 &\times\left\{
 \frac{\lambda^{-u}\epsilon^{-\alpha u}\Gamma(R+j+v)}{(1-u)_j}
 -\frac{\mu^{-u}\epsilon^{-\beta u}
              \Gamma(1-u)\Gamma(R+j+u+v)}{j!}\right\}.
\end{align}
For each fixed $L$, this is holomorphic in $(u,v)$.
Its constant in $L$ after any fixed spectral differentiation is
obtained by setting $L=0$ before differentiating; the result is
\eqref{mc:singular-contact}.

Choose a finite integer $J$ with $(J+1)(\alpha-\beta)>\beta R$,
then choose the spectral radius $\eta>0$ small enough that also
$(J+1)(\alpha-\beta)-\beta R-(\alpha+\beta)\eta>0$.
On this sufficiently small spectral disk the tail $j>J$ in the
paired series tends uniformly to zero: the coefficient growth
is at most polynomial in $j$, while $\rho^j$ is geometric.
The divided difference in $u$ introduces only a fixed logarithmic
factor, which can be absorbed by shrinking the spectral disk.
Thus only $E_j=0$ can contribute a constant, and this is exactly
the condition defining $\mathcal H_*$. Together with the preceding
mixed-term analysis this proves Theorem~\ref{mc:main} and also
provides all divergent terms through a finite truncation of
\eqref{mc:uniform-expansion} and \eqref{mc:two-branches}.

\subsection{Harmonic and zeta evaluations of the new contacts}

At spectral degree zero, the singular contact is particularly
simple. With $H_0=0$,
\begin{equation}\label{mc:H00}
 \boxed{\displaystyle
 \mathcal H_j(0,0)=(-1)^t\frac{(R+j-1)!}{j!}
       \frac{\lambda^j}{\mu^{R+j}}
             \left(\log\frac\mu\lambda+H_j-H_{R+j-1}\right).}
\end{equation}
This follows by differentiating the numerator of
\eqref{mc:singular-contact} once in $u$.
The derivative of $\log(1-u)_j^{-1}$ is $H_j$ at zero;
the derivative of
$\log\Gamma(1-u)+\log\Gamma(R+j+u)$ is $H_{R+j-1}$.
Formula \eqref{mc:H00} recovers the second resonance family in
Corollary~\ref{mc:ordered-CT}.

For arbitrary spectral degrees, the finite Taylor computations use
\begin{align}\label{mc:harmonic-jets}
 \log\frac{j!}{(1-u)_j}&=\sum_{h\ge1}\frac{H_j^{(h)}}h u^h,\notag\\
 \log\Gamma(N+u)-\log\Gamma(N)
 &=\psi(N)u+
       \sum_{h\ge2}\frac{(-1)^h}{h}
                 \bigl(\zeta(h)-H_{N-1}^{(h)}\bigr)u^h.
\end{align}
Truncate these convergent expansions at the required degree and
exponentiate, equivalently using complete Bell polynomials.
This proves the stated finite constant alphabet constructively.

For completeness, the singular contact can be evaluated without
leaving even a Gamma derivative in the answer. Put $M=R+j$ and
$D=m+n+1$, and form two finite polynomials:
\begin{align}\label{mc:finite-logarithms}
 Q_D(u,v)={}&-u\log\lambda-v\log\mu
       +\sum_{h=1}^D\frac{H_j^{(h)}}h u^h
       +\sum_{h=1}^D\frac{(-1)^{h+1}H_{M-1}^{(h)}}h v^h,
       \notag\\
 T_D(u,v)={}&-(u+v)\log\mu
       +\sum_{h=1}^D\frac{(-1)^{h+1}H_{M-1}^{(h)}}h(u+v)^h
       \notag\\
 &+\sum_{h=2}^D\frac{\zeta(h)}h
          \left[u^h+(-1)^h\bigl((u+v)^h-v^h\bigr)\right].
\end{align}
Define the finite exponential polynomial
$\mathcal E_D(X)=\sum_{k=0}^D X^k/k!$.
Then the evaluated contact is
\begin{equation}\label{mc:finite-jet}
\begin{split}
 &(-1)^{m+n}\partial_u^m\partial_v^n\mathcal H_j(0,0)\\
 &\quad=
 \frac{(-1)^{m+n+t}m!n!(M-1)!}{j!}
     \frac{\lambda^j}{\mu^M}
 [u^{m+1}v^n]\left\{\mathcal E_D(Q_D)-\mathcal E_D(T_D)\right\}.
\end{split}
\end{equation}
The notation $[u^av^b]$ extracts a coefficient from an ordinary
finite polynomial. To prove the formula, divide the two Gamma products
in \eqref{mc:singular-contact} by their common constant $(M-1)!/j!$.
Their logarithms, including the amplitude factors, have Taylor
polynomials $Q_D$ and $T_D$. Euler's constant cancels in each.
Both polynomials have zero constant term, so terms beyond the $D$th
power of their exponential cannot affect the requested coefficient.
Division by the removable $u$ is exactly extraction at $u^{m+1}$.
This gives a finite expression in harmonic numbers and ordinary zeta
values at every specified pair of Stieltjes indices. The two mixed
contacts are obtained just as finitely by multiplying the polynomials
$(1+u)_k$, the exponential of $-u\log\lambda$, and the zeta jets
in \eqref{mc:g-explicit}.

There is an especially transparent first-Stieltjes example. Take
$r=t=0$, so $R=1$. A singular contact occurs exactly when
$\alpha/\beta=1+1/j$, $j\ge1$. Direct expansion of
\eqref{mc:singular-contact} yields
\begin{align}\label{mc:H10}
 -\partial_u\mathcal H_j(0,0)
 =\frac{\lambda^j}{\mu^{j+1}}\bigg[&\zeta(2)-H_j^{(2)}\notag\\
 &+\frac{(H_j-\log\mu)^2-(H_j-\log\lambda)^2}{2}\bigg].
\end{align}
Consequently, for $a=\epsilon^{j+1}$ and $b=\epsilon^j$,
\begin{equation}\label{mc:stieltjes-tail-contact}
 \boxed{\displaystyle
 \operatorname{CT}_{\epsilon}\int_0^1
          \gamma_1(x+\epsilon^{j+1})\gamma_0(x+\epsilon^j)dx
   =\frac32\gamma_2+2\gamma\zeta(2)-\zeta(3)
                         +\zeta(2)-H_j^{(2)}.}
\end{equation}
The new term is the ordinary trigamma tail
$\zeta(2)-H_j^{(2)}=\psi'(j+1)$.
Thus all positive integer trigamma tails appear as exact
Stieltjes collision contacts.

This family also has a classical polylogarithmic generator:
for $|z|<1$,
\begin{equation}\label{mc:polylog-generator}
 \sum_{j\ge1}\bigl(\zeta(2)-H_j^{(2)}\bigr)z^j
       =\frac{\zeta(2)z-\operatorname{Li}_2(z)}{1-z}.
\end{equation}
Its ordinary primitive can be normalized at zero using
\begin{equation}\label{mc:polylog-primitive}
 \int_0^z\frac1t
       \sum_{j\ge1}\bigl(\zeta(2)-H_j^{(2)}\bigr)t^jdt
   =2\operatorname{Li}_{2,1}(z)-\operatorname{Li}_3(z)
      +(\operatorname{Li}_2(z)-\zeta(2))\log(1-z).
\end{equation}
Here $\operatorname{Li}_{2,1}(z)=
\sum_{n>k\ge1}z^n/(n^2k)$, on its analytic branch at zero.
The generating identity is an elementary interchange of absolutely
convergent sums. The primitive identity is checked by differentiation
using $\frac{d}{dz}\operatorname{Li}_{2,1}(z)
=\operatorname{Li}_{1,1}(z)/z=\log^2(1-z)/(2z)$ together with
the corresponding product derivatives; an independent quadrature check
is recorded with the computational artifacts. In fact classical
polylogarithms suffice: for $0<z<1$, the same primitive equals
\begin{equation}\label{mc:classical-primitive}
 2\zeta(3)-\operatorname{Li}_3(z)-2\operatorname{Li}_3(1-z)
                +\log(1-z)\operatorname{Li}_2(1-z).
\end{equation}
To verify this expression directly, differentiate and use Euler's
dilogarithm reflection identity; its limit at zero is zero. At
$z=1$ it gives the endpoint identity
$\sum_{j\ge1}(\zeta(2)-H_j^{(2)})/j=\zeta(3)$.
The elementary generating and primitive identities are classical
Euler-sum consequences; their role here is to evaluate the newly
identified contact family, not to claim those consequences as
independently new polylogarithm identities.

\subsection{What the different-scale theorem changes}

The fixed-rate source remains correct with its stated hypotheses.
It must not be extrapolated by substituting an
$\epsilon$-dependent rate in a remainder proved only for fixed rates.
Two changes are essential: high Taylor coefficients at one can
survive when $b^\ell/a^{r-k}$ has exponent zero, and the two
singular Frobenius branches can jointly have exponent zero.
Both effects are present in \eqref{mc:ordered-examples}; the
first-index law \eqref{mc:stieltjes-tail-contact} isolates the
second effect without argument derivatives.

The theorem does not assert uniformity as the exponent ratio
approaches a resonance while $\epsilon$ also varies. In that
transition, a factor $\epsilon^{E_j}$ must be retained rather than
replaced by its constant coefficient. Formula \eqref{mc:H-log}
provides the exact local term for such a further analysis.

% END sections/03_multiscale_collisions.tex

% BEGIN sections/04_barnes_companion.tex
% Research notes for the parent report.  This is an input fragment.
% Standard amsmath, amssymb, amsthm theorem/proof environments are assumed.

\section{A Barnes-integral companion}
\label{bim:section}

This section records evaluated consequences of the classical first Barnes
lemma.  The underlying integral is not new: it is NIST DLMF
\href{https://dlmf.nist.gov/5.13.E3}{Equation 5.13.3} \cite{DLMF}.  The present
contribution is the selection and complete differentiation of the
centered fourth moment and the pure trigamma correlation.  No historical priority claim is made for these specializations.  These ordinary integrals are
distinct from the higher-dimensional Gamma density used in the recent
\emph{Gamma Normalizations and Colored Moments} report.

For complex parameters of positive real part, put
\begin{align}
 W(a,b,c,d;x)&=\Gamma(a+ix)\Gamma(b+ix)
                  \Gamma(c-ix)\Gamma(d-ix),\\
 F(a,b,c,d)&=\frac{\Gamma(a+c)\Gamma(a+d)
                       \Gamma(b+c)\Gamma(b+d)}
                      {\Gamma(a+b+c+d)}.
\end{align}
Barnes' first lemma states
\begin{equation}\label{bim:barnes}
 \frac1{2\pi}\int_{-\infty}^{\infty}W(a,b,c,d;x)\,dx
                              =F(a,b,c,d).
\end{equation}
Every finite parameter derivative can be passed under this integral.
Indeed, on a compact parameter set inside the four right half-planes,
Stirling's formula gives exponential decay $e^{-2\pi|x|}$ times a
fixed polynomial in $1+|x|$; differentiating adds only a fixed power of
$1+\log(2+|x|)$.  This also proves that all integrals below are ordinary
absolutely convergent integrals.

\subsection{Centered mixed moments}

For $\sigma>0$ define
\begin{equation}\label{bim:density}
 p_\sigma(x)=\frac{\Gamma(4\sigma)}
                   {2\pi\Gamma(2\sigma)^4}
                        |\Gamma(\sigma+ix)|^4,
 \qquad \int_{\mathbb R}p_\sigma(x)\,dx=1.
\end{equation}
Write $\mathbb E_\sigma$ for integration against this density, and put
\begin{equation}\label{bim:center}
 m_\sigma=2\psi(2\sigma)-\psi(4\sigma),\qquad
 U_\sigma(x)=\psi(\sigma+ix)-m_\sigma.
\end{equation}
The following formulas involve only ordinary polygamma values:
\begin{align}
 \mathbb E_\sigma U_\sigma&=0,\label{bim:mean}\\
 \mathbb E_\sigma U_\sigma^2&=-\psi'(4\sigma),
                                               \label{bim:square}\\
 \mathbb E_\sigma|U_\sigma|^2
          &=\psi'(2\sigma)-\psi'(4\sigma),      \label{bim:abs-square}\\
 \mathbb E_\sigma\bigl(U_\sigma|U_\sigma|^2\bigr)
          &=-\psi''(4\sigma),                  \label{bim:cubic}\\
 \mathbb E_\sigma|U_\sigma|^4
          &=\psi'(4\sigma)^2
             +2\bigl(\psi'(2\sigma)-\psi'(4\sigma)\bigr)^2
             -\psi'''(4\sigma).                \label{bim:fourth}
\end{align}
The left side of \eqref{bim:cubic} is real by the change $x\mapsto-x$.

\begin{proof}
Set $L=\log F$.  At $a=b=c=d=\sigma$, all first logarithmic
derivatives equal $m_\sigma$.  Subtract this linear term by replacing
$F/F(\sigma,\sigma,\sigma,\sigma)$ with
\[
 \widetilde F(a,b,c,d)
 =e^{-m_\sigma(a+b+c+d-4\sigma)}
      \frac{F(a,b,c,d)}{F(\sigma,\sigma,\sigma,\sigma)}.
\]
Its single-slot integral insertions are $U_\sigma$ in the slots $a,b$
and $\overline U_\sigma$ in $c,d$.  The first derivatives of
$\log\widetilde F$ vanish.  For distinct slots the remaining derivatives
needed here are
\begin{align*}
 L_{ab}=L_{cd}&=q=-\psi'(4\sigma),\\
 L_{ac}=L_{ad}=L_{bc}=L_{bd}
      &=p=\psi'(2\sigma)-\psi'(4\sigma),\\
 L_{abc}=L_{abd}=L_{acd}=L_{bcd}&=-\psi''(4\sigma),\\
 L_{abcd}&=-\psi'''(4\sigma).
\end{align*}
The derivatives of $\widetilde F=e^{\log\widetilde F}$ of orders two
and three give the first four assertions.  Its $abcd$ derivative is
$L_{ab}L_{cd}+L_{ac}L_{bd}+L_{ad}L_{bc}+L_{abcd}$, because every
partition containing a singleton has zero contribution.  This is
$q^2+2p^2-\psi'''(4\sigma)$ and proves \eqref{bim:fourth}.
\end{proof}

\subsection{Eliminating all digamma insertions from a trigamma correlation}

\begin{theorem}\label{bim:trigamma}
For $\sigma>0$ one has
\begin{align}
 \mathbb E_\sigma\psi'(\sigma+ix)&=2\psi'(2\sigma),
                                                   \label{bim:trigamma-mean}\\
 \mathbb E_\sigma\left|\psi'(\sigma+ix)\right|^2
                  &=4\psi'(2\sigma)^2+\psi'''(2\sigma).
                                                   \label{bim:trigamma-square}
\end{align}
More generally, if $\Re u,\Re v>0$ and $s=u+v$, then
\begin{equation}\label{bim:unequal-trigamma}
\begin{split}
 &\int_{\mathbb R}\Gamma(u+ix)^2\Gamma(v-ix)^2
                         \psi'(u+ix)\psi'(v-ix)\,dx\\
 &\qquad=2\pi\frac{\Gamma(s)^4}{\Gamma(2s)}
                           \bigl(4\psi'(s)^2+\psi'''(s)\bigr).
\end{split}
\end{equation}
\end{theorem}

\begin{proof}
Use the constant-coefficient operators
\[
 D_+=\partial_a^2-\partial_a\partial_b,
 \qquad D_-=\partial_c^2-\partial_c\partial_d.
\]
After setting $a=b=u$ the first operator inserts exactly
$\psi'(u+ix)$, since the two copies of $\psi(u+ix)^2$ cancel.  The
second operator inserts $\psi'(v-ix)$ after $c=d=v$.

For an explicit right-side calculation, center as above, now at
$a=b=u,c=d=v$; the common first logarithmic derivative is
$2\psi(s)-\psi(2s)$.  Put
\[
 \begin{gathered}
 A=\psi'(s),\quad B=\psi'(2s),\quad
 A_3=\psi'''(s),\quad B_3=\psi'''(2s),\\
 d=2A-B,\quad q=-B,\quad p=A-B.
 \end{gathered}
\]
Centering does not affect $D_+D_-\widetilde F$ at coalescence: the extra
terms contain $\partial_a-\partial_b$ or
$\partial_c-\partial_d$, which vanish by the respective symmetries.
The required derivatives are
\begin{align*}
 \widetilde F_{aacc}&=d^2+2p^2+A_3-B_3,\\
 \widetilde F_{aacd}=\widetilde F_{abcc}&=dq+2p^2-B_3,\\
 \widetilde F_{abcd}&=q^2+2p^2-B_3.
\end{align*}
Consequently
\[
 D_+D_-\widetilde F=(d-q)^2+A_3=4A^2+A_3.
\]
Multiplying by $2\pi F(u,u,v,v)$ proves
\eqref{bim:unequal-trigamma}; setting $u=v=\sigma$ proves
\eqref{bim:trigamma-square}.  The second-order calculation
$\widetilde F_{aa}-\widetilde F_{ab}=d-q=2A$ gives
\eqref{bim:trigamma-mean}.
\end{proof}

The unequal-shift formula can equally be written using
$\psi'(z)=\zeta(2,z)$ and $\psi'''(z)=6\zeta(4,z)$.  Thus it is a
Gamma-weighted Hurwitz-zeta convolution with two independently chosen
argument shifts.  The only shift invariant on the right is their sum.

\subsection{Explicit half- and quarter-argument evaluations}

At $\sigma=1/2$, the Gamma reflection formula makes
$p_{1/2}(x)=(\pi/2)\operatorname{sech}^2(\pi x)$ and
$U_{1/2}(x)=\psi(1/2+ix)+\gamma+1$.  Equations
\eqref{bim:abs-square}, \eqref{bim:fourth}, and
\eqref{bim:trigamma-square} become
\begin{align}
 \int_{\mathbb R}\frac{|\psi(1/2+ix)+\gamma+1|^2}
                            {\cosh^2(\pi x)}\,dx
                   &=\frac2\pi,\label{bim:half-square}\\
 \int_{\mathbb R}\frac{|\psi(1/2+ix)+\gamma+1|^4}
                            {\cosh^2(\pi x)}\,dx
                   &=\frac2\pi\left(9-\frac{\pi^2}{3}
                                          -\frac{7\pi^4}{180}\right),
                                               \label{bim:half-fourth}\\
 \int_{\mathbb R}\frac{|\psi'(1/2+ix)|^2}
                            {\cosh^2(\pi x)}\,dx
                   &=\frac{16\pi^3}{45}.\label{bim:half-trigamma}
\end{align}
Here the cancellation of Euler's constant is exact after centering.

At $\sigma=1/4$, put $V(x)=\psi(1/4+ix)+\gamma+4\log2$.
The normalization is $p_{1/4}(x)=|\Gamma(1/4+ix)|^4/(2\pi^3)$, and
one obtains the equally short evaluations
\begin{align}
 \int_{\mathbb R}|\Gamma(1/4+ix)|^4|V(x)|^2\,dx&=\frac{2\pi^5}{3},
                                                   \label{bim:quarter-square}\\
 \int_{\mathbb R}|\Gamma(1/4+ix)|^4|V(x)|^4\,dx&=\frac{11\pi^7}{30},
                                                   \label{bim:quarter-fourth}\\
 \int_{\mathbb R}|\Gamma(1/4+ix)|^4|\psi'(1/4+ix)|^2\,dx&=4\pi^7.
                                                   \label{bim:quarter-trigamma}
\end{align}
The special-value reductions use DLMF
\href{https://dlmf.nist.gov/5.15.E2}{5.15.2} and
\href{https://dlmf.nist.gov/5.15.E3}{5.15.3}.


% END sections/04_barnes_companion.tex

% BEGIN sections/05_source_status.tex
% Parent-ready source/status fragment. No preamble.
\section{Source audit, integration decisions, and further research}
\label{status:section}

The source anchor for this continuation is ProveIt revision
\texttt{03112a7d3b91} \cite{ProveIt}.
The audit concerns the incoming passages used in the proofs and the
manuscript's relevant status records. It is not a certification of
every result in the repository. The three main advances have the
following precise relation to that source state.

\subsection{What the present results change}

\begin{enumerate}
\item \textbf{Joint Gamma--polygamma moments.}
The Gamma-normalization report's
\path{sections/02_gaussian.tex} proves the paired deformation and
the single-trigamma Gaussian moments through degree $n+1$;
research question Q2 in
\path{sections/06_audit_research.tex} asks for the joint
polygamma moments \cite{GammaIncoming}.
Theorem~\ref{gc:main} separates all joint cumulants through total
weight $2n+2$, with weight $q+1$ assigned to $R_q$.
It includes digamma and even polygamma orders. Formula
\eqref{gc:moment} then evaluates every retained ordinary moment,
and Corollary~\ref{gc:wick} gives the centered positive-order Wick
rule. This answers the joint-separation question in the stated
weighted range; the beyond-range classification remains open.
The nonzero small-$\sigma$ skewness limit in
Proposition~\ref{gc:sharpness} supplies an obstruction relevant to
the same report's Q1, but does not evaluate the general first
excluded trigamma moment.
The underlying Gustafson integral remains classical
\cite{DerkachovManashov}; its Gamma normalizer and the paired law
retain their incoming attribution.

\item \textbf{Effective reconstruction and the complete gap-eight ray.}
Section~8 of the ordered-boundary report gives eventual
nonsingularity at each fixed gap; Section~9, research question~1,
asks for an effective threshold and finite completion
\cite{OrderedIncoming}. Theorem~\ref{eg:effective} supplies an
explicit computable $B(m)$ for every fixed half-gap $m$.
The later gap-six report's
\path{sections/05_audit_and_questions.tex}, question Q1,
explicitly asks for complete gap-eight coverage \cite{GapSix}.
Theorem~\ref{eg:eight-main} proves this for every
$(p,q)=(2b-8,2b)$, $b\ge5$, using a proved analytic tail and an
exact finite-field completion. Corollary~\ref{eg:eight-identities}
turns that result into ordinary convergent polygamma-product
identities. The unrestricted all-pairs prefix conjecture remains
open. The finite certificate proves nonvanishing, and does not
prove positivity at every intervening finite index.

\item \textbf{Stieltjes collisions at different monomial scales.}
The gap-six report's Section~4 treats fixed positive rates
$a=\lambda\epsilon$, $b=\mu\epsilon$; its question Q10 asks for
$a=\lambda\epsilon^\alpha$, $b=\mu\epsilon^\beta$ with
$\alpha\ne\beta$ \cite{GapSix}.
Theorem~\ref{mc:main} evaluates the constant for every positive
monomial pair and every fixed pair of Stieltjes indices and
argument derivatives. Its proof also specifies all divergent
terms through that constant by finite truncation. The three
resonance conditions in \eqref{mc:three-mechanisms} identify the
additional contacts, including the singular--singular contact.
The earlier lattice-contact report's Q9 asked only for fixed
unequal rates \cite{LatticeIncoming}; that problem was already
answered by the gap-six report and is not reclaimed here.
General nonmonomial paths and three independently shifted factors
remain outside the present theorem. The trigamma-tail contact
family \eqref{mc:stieltjes-tail-contact} and its polylogarithmic
generator and primitive are explicit evaluated consequences of
the new scale classification.
\end{enumerate}

The Barnes-integral companion supplies additional ordinary integral
evaluations obtained by explicitly differentiating a classical Barnes
lemma \cite{DLMF}. It should be credited in that form. A bounded
literature search and comparison with the inspected incoming material
support the stated continuation choices; they do not establish global
historical novelty.

\subsection{Errors, status corrections, and essential qualifications}

No newly false inherited mathematical theorem was found in the
specific dependencies used for these results. The changes above are
proved extensions and research-status updates. They should not be
presented as corrections to formulas which the sources stated correctly.
In particular, the incoming Gamma report already restricts its finite
Gaussian law, and the gap-six report already distinguishes finite-field
nonvanishing from a real-sign theorem.

One source-list reconciliation is necessary. The gap-six report's Q4
and the ordered-boundary report's question~2 ask for evaluations of
$K_m$. The later Gamma-normalization report already supplies them
\cite{GammaIncoming,GapSix,OrderedIncoming}. These questions should be
marked answered by that predecessor, rather than counted among the
present results. The repeated README paragraph about the 1,032-input
first-letter-$X$ family, rank 1,008 and kernel dimension 24 remains
present at the pinned revision. Earlier incoming reports already
identified this editorial duplication. Removing its second occurrence
would be an inherited editorial cleanup, not a new mathematical erratum.

Three conventions are material to the proofs. The Gamma moment
classification retains its total-weight restriction: beyond that range,
Newton relations prevent unrestricted comparison of power-sum
coefficients. The Bernoulli reconstruction matrix retains the omission
of a lower binomial index zero; evaluating that entry as
$\binom n0=1$ changes the problem. The Stieltjes constant is taken in
the displayed small parameter, with the fixed-coordinate finite part
kept distinct from the path-dependent constant. Spectral pole pairs
must be combined before their removable values and derivatives are
taken. These are integration requirements, not newly discovered
mistakes in the cited sources.

The manuscript's short $S_6$ and revised $S_8$ candidates remain open;
the old rejected $S_8$ candidate remains rejected. None of the Gamma
moment, reconstruction, or collision theorems implies those polylogarithmic
period identities. Numerical proximity and rank in a specified formal
relation space do not supply their missing exact equalities.
The verification artifacts here likewise distinguish exact finite
certificates from floating-point diagnostics. The analytic arguments
justify convergence, continuation and coefficient extraction.

\subsection{Further questions directed toward exact identities}

\begin{enumerate}
\item \textbf{Evaluate the first excluded polygamma cumulants.}
At total weight $2n+3$ and above, determine the relations induced
by $e_{2n+3},e_{2n+4},\ldots$ and identify additional integral
identities that separate the remaining cumulants. The immediate
scalar target is the exact $(n+2)$nd trigamma Gaussian defect at
general $n,\sigma$. The proved $n=1$ small-shift limit is a
normalization constraint on any proposed formula.

\item \textbf{Develop joint moment formulas for unequal Gamma clusters.}
Allow several distinct base shifts with specified multiplicities.
Determine which mixed polygamma moments can be separated by
independent deformations within each cluster, how their finite
coordinate relations change when clusters merge, and which new
ordinary integrals arise at the first unresolved degrees. A useful
answer would give explicit polygamma or finite hypergeometric
coefficients, rather than only a differentiated generating integral.

\item \textbf{Find closed reconstruction rows and identities in $b$.}
The gap-eight theorem specifies $P_b$ by a rational inverse matrix.
Find finite-sum descriptions or recurrences for its coefficients and
for the corresponding subtraction polynomial. Decide which inverse
rows yield especially short zeta-valued identities. The exact
leading coefficient and the ordinary sum in
\eqref{eg:eight-sum} provide tests for every proposed recurrence.

\item \textbf{Replace fixed-gap finite completion by structural algebra.}
Find a recurrence invariant, factorization, or appropriate
variation-diminishing statement that proves nonvanishing for an
entire gap ray without checking every exceptional index separately.
Such a theorem must respect the known early sign behavior.
Connecting this structure with the fixed-smaller-order results is a
concrete route toward the unrestricted unequal-even prefix problem.

\item \textbf{Compute collision constants for nonmonomial paths.}
Take $a=\lambda\epsilon^\alpha A(\epsilon)$ and
$b=\mu\epsilon^\beta B(\epsilon)$, with prescribed analytic or
power--logarithm expansions and nonzero leading coefficients.
Determine the finite jets of $A,B$ that can affect a specified
constant and derive its exact transformation law under a change of
the small parameter. The monomial resonance rules identify the
base exponents with which these additional coefficients interact.

\item \textbf{Evaluate three-factor Stieltjes collision contacts.}
For three independently shifted derivatives, classify contributions
from every singular subset and pair all intersecting spectral poles
before differentiating. Seek a finite evaluated contact formula in
logarithms, generalized harmonic numbers and zeta derivatives, and
identify precisely where multivariable polylogarithms or other new
functions become necessary.

\item \textbf{Generate and integrate the higher contact families.}
For fixed Stieltjes and derivative orders, sum the integer-resonance
contacts against $z^j$ and determine their polylogarithmic alphabet.
The trigamma-tail generator and its anchored primitive are the first
case. Compute further Euler primitives, ordinary or finite-part
endpoint values, and character-weighted versions; determine which
reduce to zeta products and which retain genuine multiple
polylogarithms.

\item \textbf{Construct an explicit functional route to $S_6$ or revised
$S_8$.}
If one of the new generators or integral deformations is to inform
these period problems, first exhibit its exact map to the canonical
residual, including all products and branch constants. A new
functional or iterated-integral identity outside the tested relation
spaces would constitute progress; an unrelated rank theorem or
additional decimal agreement would not.
\end{enumerate}

% END sections/05_source_status.tex

% BEGIN sections/06_reproducibility.tex
\section{Reproducibility and the role of computation}
\label{repro:section}

The package contains the complete article as one standalone TeX file,
the modular source, its compiled PDF, all verification programs, and
the recorded results. The exact source revision and the inherited
inputs are identified in the bibliography. No running service,
unpublished data, or repository mutation is required to reproduce the
calculations. The mathematical assertions are separated into three
types of evidence.

\paragraph{Analytic proofs.}
The joint-cumulant theorem is proved by a classical Gamma integral and
finite polynomial injectivity. The effective reconstruction tail is
proved by an explicit two-term phase expansion and a rational
determinant margin. The monomial collision theorem is proved by
uniform local subtraction, finite selection of exponents, and
holomorphic cancellation of paired singular branches. The ordinary
Barnes evaluations follow by differentiated integral identities.
The limit and sum convergence arguments are part of those proofs.

\paragraph{Exact finite computation.}
The C++ program \texttt{verify\_gap\_eight.cpp} uses only unsigned
integer arithmetic and arithmetic in the prime field
$\mathbb F_{2147483647}$. All $8187$ determinant residues for
$5\le b\le8191$ are explicitly recorded. The verifier checks the
prime and every needed denominator and independently checks all
$8194$ cotangent coefficients against the sine and cosine series.
An independent replay of the compiled program reproduced the complete
certificate exactly. The SymPy bridge program separately checks the
original Bernoulli matrices and their rational determinant reduction
for $5\le b\le16$, and the concrete $(14,22)$ sum has its own exact
generator. The finite-field certificate proves rational nonvanishing;
it does not estimate a real determinant or deduce its sign.

\paragraph{Independent symbolic and numerical diagnostics.}
The Gamma program performs $122$ exact algebra assertions and $17$
integral checks. The latter extract cumulants from determinants of
independently integrated moment entries; they do not substitute the
same Gamma quotient on both sides. The two precision runs and their
residuals are described in Section~\ref{gc:section}.
The Stieltjes compiler implements the finite contact rule using exact
SymPy expressions. Its test program performs $72$ exact assertions,
including comparisons with derivatives of the elementary logarithmic
kernel, the first-index dilogarithmic contact, and the inherited
fixed-rate law. Five nonreal spectral checks compare the paired branch
series with the original convergent kernel integral; three further
checks compare the classical primitive with direct integration.
At $64$ digits the largest scaled spectral residual is below
$8.2\times10^{-56}$.

The separate \texttt{identity\_audit.py} evaluates ordered collision
remainders and six Barnes moment formulas at $90$ decimal digits.
The Barnes residuals are below $3\times10^{-91}$. The collision
diagnostics deliberately include a nonresonant irrational exponent
ratio, for which a small positive exponent makes convergence slow.
That example is recorded with its vanishing-term acceleration rather
than mislabeled as a rapidly convergent check. All floating-point
residuals are diagnostics, not rigorous numerical error intervals.
The exact theorem does not depend on their size.

\paragraph{How to reproduce.}
The \texttt{README.md} gives the precise build and verification
commands. The minimal numerical dependencies are Python~3,
SymPy, and mpmath; the certificate additionally needs a C++17
compiler. The TeX uses standard mathematical packages and compiles
with pdfLaTeX through \texttt{latexmk}. Internal analytical audit
notes identify the delicate normalizations and remainder arguments.
These are separate derivations within this research workflow and are
not a claim of external peer review or proof-assistant certification.

% END sections/06_reproducibility.tex

% BEGIN references.tex
\begin{thebibliography}{99}

\bibitem{ProveIt}
V.~Reshetnikov, \emph{ProveIt}, polylogarithm manuscript and incoming
research reports, repository revision
\texttt{03112a7d3b91efdd3c00b6e595cf0ffa4694d2dc},
11 October 2026 (repository commit date), inspected 12 October 2026.
\href{https://github.com/VladimirReshetnikov/ProveIt/tree/03112a7d3b91efdd3c00b6e595cf0ffa4694d2dc/Analysis/Polylogarithms/docs/manuscript}{Pinned manuscript};
\href{https://github.com/VladimirReshetnikov/ProveIt/tree/03112a7d3b91efdd3c00b6e595cf0ffa4694d2dc/docs/incoming}{pinned incoming directory}.

\bibitem{OrderedIncoming}
ProveIt incoming report, \emph{Ordered Boundary Identities},
11 October 2026, archive
\path{ProveIt_Ordered_Boundary_Identities_2026-10-11.zip}
in \cite{ProveIt}. The inherited fixed-gap recurrence, Gamma basis,
complementary-minor bridge, and eventual nonsingularity theorem.

\bibitem{GapSix}
ProveIt incoming report, \emph{Gap-Six Harmonic Determinants},
12 October 2026, archive
\path{ProveIt_Gap_Six_Harmonic_Determinants_2026-10-12.zip}
in \cite{ProveIt}. See \path{sections/01_gap_six.tex},
\path{sections/04_collisions.tex}, and \path{sections/05_audit_and_questions.tex}.

\bibitem{GammaIncoming}
ProveIt incoming report, \emph{Gamma Normalizations and Colored Moments},
12 October 2026, archive
\path{ProveIt_Gamma_Normalizations_and_Colored_Moments_2026-10-12.zip}
in \cite{ProveIt}. See \path{sections/02_gaussian.tex} and
\path{sections/06_audit_research.tex}, especially research question~2.

\bibitem{LatticeIncoming}
ProveIt incoming report, \emph{Lattice Contacts and Mixed Zeta Identities},
11 October 2026, archive
\path{ProveIt_Lattice_Contacts_and_Mixed_Zeta_Identities_2026-10-11.zip}
in \cite{ProveIt}.

\bibitem{DerkachovManashov}
S.~E. Derkachov and A.~N. Manashov,
\emph{Spin Chains and Gustafson's Integrals},
Journal of Physics A: Mathematical and Theoretical \textbf{50} (2017),
294006. \href{https://doi.org/10.1088/1751-8121/aa749a}{doi:10.1088/1751-8121/aa749a};
\href{https://arxiv.org/abs/1611.09593v2}{arXiv:1611.09593v2}.
The symplectic Barnes integral used here is equation~(2).

\bibitem{DLMF}
F.~W.~J. Olver et al., eds.,
\emph{NIST Digital Library of Mathematical Functions},
\url{https://dlmf.nist.gov/}, accessed 12 October 2026.
In particular \href{https://dlmf.nist.gov/5.11}{\S5.11}
(Gamma asymptotics and remainders),
\href{https://dlmf.nist.gov/5.13}{\S5.13}
(Barnes integrals), \href{https://dlmf.nist.gov/5.15}{\S5.15}
(polygamma functions), \href{https://dlmf.nist.gov/15.8}{\S15.8}
(hypergeometric transformations), and
\href{https://dlmf.nist.gov/25.11}{\S25.11}
(Hurwitz zeta functions).

\bibitem{TicaGalapon}
C.~D. Tica and E.~A. Galapon,
\emph{Finite-part integration of the generalized Stieltjes transform
and its dominant asymptotic behavior for small values of the parameter.
II. Non-integer orders}, Journal of Mathematical Physics
\textbf{60} (2019), 013502.
\href{https://doi.org/10.1063/1.5038900}{doi:10.1063/1.5038900}.

\bibitem{ShpotParis}
M.~A. Shpot and R.~B. Paris,
\emph{Integrals of products of Hurwitz zeta functions via Feynman
parametrization and two double sums of Riemann zeta functions},
\href{https://arxiv.org/abs/1609.05658v2}{arXiv:1609.05658v2}
(2016; revised 2020).

\end{thebibliography}

% END references.tex

\end{document}
